%% POPL 2026 submission
\documentclass[acmsmall,sigplan,anonymous,review,screen]{acmart}

\acmConference[POPL '26]{Symposium on Principles of Programming Languages}{January 2026}{}
\acmBooktitle{Proceedings of the 53rd Annual ACM Symposium on Principles of
  Programming Languages (POPL '26)}
\acmDOI{}
\acmISBN{}

%% --- Packages ---
\usepackage{amsmath,amsthm}
\usepackage{mathtools}
\usepackage{xspace}
\usepackage{microtype}
\usepackage{booktabs}
\usepackage{listings}
\usepackage{enumitem}
\usepackage{xcolor}
\usepackage{array}
\usepackage{bussproofs}

% Simple inference rule macro
\newcommand{\inferrule}[2]{%
  \dfrac{\displaystyle #1}{\displaystyle #2}%
}

\theoremstyle{plain}
\newtheorem{theorem}{Theorem}[section]
\newtheorem{lemma}[theorem]{Lemma}
\newtheorem{corollary}[theorem]{Corollary}
\newtheorem{proposition}[theorem]{Proposition}

\theoremstyle{definition}
\newtheorem{definition}[theorem]{Definition}
\newtheorem{example}[theorem]{Example}
\newtheorem{remark}[theorem]{Remark}

%% --- Macros ---
\newcommand{\qSL}{\textsc{qSL}\xspace}
\newcommand{\CTIMP}{\textsc{ct-imp}\xspace}
\newcommand{\emerge}{c}
\newcommand{\rs}{\mathit{rs}}
\newcommand{\fp}{\mathit{fp}}
\newcommand{\emp}{\mathit{emp}}
\newcommand{\sep}{*}
\newcommand{\wand}{-\!\!*}
\newcommand{\HS}{H\!H}
\newcommand{\HH}[2]{H\!H^{#1}(#2)}
\newcommand{\norm}[1]{\|#1\|}
\newcommand{\defect}{\varepsilon}
\newcommand{\lean}[1]{\texttt{#1}}
\newcommand{\leancite}[2]{\lean{#1}\,[\texttt{#2}]}
\newcommand{\compfootprint}{\mathit{fp}}

\lstset{
  basicstyle=\footnotesize\ttfamily,
  columns=flexible,
  keepspaces=true,
  mathescape=true,
}

\setcopyright{none}

%% =================================================================
\begin{document}

\title{Cocycle-Graded Separation Logic: Decidable Equivalence Modulo
       Bounded Compositional Defect}

\begin{abstract}
Reynolds's frame rule asserts that the framing resource $R$ is unaffected
by a command $c$; the rule is exact when resource interactions are additive.
In practice, cache eviction, speculative execution, and probabilistic
scheduling all produce measurable, bounded interaction between the program
footprint and the frame resource.  We give a uniform algebraic account of
this phenomenon.

Given a separation algebra $(I, \oplus, e)$ equipped with an
\emph{interference measure} $\rs : I \times I \to \mathbb{R}$ additive
in its second argument, the \emph{interference cochain}
$c(a, b, d) := \rs(a \oplus b, d) - \rs(a, d) - \rs(b, d)$
is always a Hochschild 2-cocycle of the resource monoid.  We prove this
identity fully in Lean~4 (\lean{cocycleClosed}, 9 lines, zero
\lean{sorry}).

We embed this into \qSL, a quantitative separation logic whose
\emph{cocycle-graded frame rule} accumulates $\norm{c}$ as a defect bound
and whose zero-defect case is exactly Reynolds--O'Hearn separation logic.
The logic unifies three previously separate formalisms: classical SL
($c \equiv 0$), graded effect systems (grade annotation on the resource
monoid), and quantitative type theory (slack bound on the defect).
Soundness is proved with respect to a step-indexed denotational model.

Our second contribution is decidability.  For finitely-generated
commutative resource monoids with integer-valued pairing of rank $r$
bounded by $B$, the $k$-cocycle equivalence problem has a finite quotient
monoid $M_k$ of size $O((2Bk{+}1)^r \cdot |M|)$, yielding a decision
procedure running in time $O(n \cdot |M|^2 \cdot (2Bk{+}1)^r)$.  With $k$
in unary the problem is in PSPACE; we prove PSPACE-hardness by a reduction
from TQBF, establishing PSPACE-completeness.  A cache-timing case study and
a 600-line Lean~4 mechanization of the algebraic core complete the
presentation.
\end{abstract}

\keywords{separation logic, Hochschild cohomology, graded effects,
  quantitative types, decidability, PSPACE-completeness, cache timing, Lean~4}

\maketitle

% ==================================================================
\section{Introduction}
\label{sec:intro}

\subsection{The Frame Rule and Its Quantitative Relatives}
\label{sec:intro:frame}

Reynolds's frame rule~\cite{reynolds2002separation} is the cornerstone of
compositional program reasoning in separation logic:
\begin{equation}
  \label{eq:frame}
  \inferrule{\{P\}\;c\;\{Q\}}{\{P \sep R\}\;c\;\{Q \sep R\}} \quad\text{[Frame]}
\end{equation}
The rule is \emph{exact}: the framing resource $R$ is neither consumed
nor modified by $c$.  Its soundness rests on the \emph{locality} axiom of
O'Hearn et al.~\cite{ohearn2001local}: commands only access memory within
their footprint, leaving disjoint resources unchanged.

Locality is categorical in the idealised heap model: a command either
touches a location or it does not.  Contemporary systems violate this
in a soft, computable way.

\begin{itemize}[leftmargin=1.5em,itemsep=2pt]
  \item \textbf{Cache timing.}
    A \texttt{load} by $c$ can evict cache lines belonging to the frame
    resource $R$, adding a timing penalty of $\Delta = t_{\mathrm{miss}} -
    t_{\mathrm{hit}}$ per shared line.  Crucially, this cost is
    \emph{sub-additive}: if $c$ already accessed a line $\ell$ from its
    own footprint, $\ell$'s eviction from $R$'s working set is not an
    extra cost.
  \item \textbf{GC interaction.}  A garbage-collection cycle triggered
    inside $c$ can move objects reachable from $R$, degrading pointer
    validity.  The GC cost is proportional to the live set, which
    includes $R$'s data.
  \item \textbf{Speculative execution.}  Speculative reads from $R$ inside
    $c$, as in Spectre~attacks, leak values across the architectural
    boundary.  The interference is proportional to the overlap of $c$'s
    branch predictor state with $R$'s address space.
  \item \textbf{Probabilistic scheduling.}  In concurrent programs,
    execution of $c$ shifts the scheduler probability distribution, causing
    probability mass to flow across thread-local heaps.  The \emph{cross
    defect} between $c_1 \| c_2$ is bounded by the covariance of their
    footprint distributions.
\end{itemize}

In each case the frame defect is:
(1) \emph{computable} from the footprints alone,
(2) \emph{bounded} by a constant depending on the resource monoid and
    interference measure, and
(3) \emph{compositional}: the defect for $c_1; c_2$ is at most the sum
    of individual defects (not the product).

\paragraph{The algebraic question.}
Why are frame defects always compositional?  Prior work (QSL, graded types)
introduces defect bounds as axioms or parameters.  We show that
compositionality is \emph{forced} by a universal algebraic identity: the
interference measure $\rs$ gives rise to a 3-variable cochain $c(a,b,d)$
that is, without any extra assumption, a Hochschild 2-cocycle of the
resource monoid.  The cocycle identity $\delta^2 c = 0$ is precisely the
equation that forces defects to accumulate additively.

\paragraph{The decidability question.}
Given two programs $P$ and $Q$ and a bound $k$, is their frame-defect
difference bounded by $k$?  This is the \emph{$k$-cocycle equivalence
problem}.  We show it is PSPACE-complete for commutative finitely-generated
resource monoids, giving the first complexity-theoretic characterization
of compositional program equivalence under bounded interference.

\begin{table}[t]
  \caption{The three pillars unified by \qSL.  Each prior formalism is a
    special case with a specific instantiation of the cocycle parameter.}
  \label{tab:unification}
  \centering\small
  \begin{tabular}{@{}p{2.0cm}p{2.0cm}p{2.0cm}l@{}}
    \toprule
    Formalism & Cocycle $c$ & Defect & Ref. \\
    \midrule
    Classical SL (heap) & $c \equiv 0$ & $0$ & \cite{reynolds2002separation} \\
    Graded effects & $c \equiv 0$ & grade $m$ & \cite{katsumata2014parametric} \\
    Prob.\ SL & $c \neq 0$ & $\varepsilon$ (param.) & \cite{batz2019quantitative} \\
    \textbf{\qSL} (ours) & $c$ (derived) & $\norm{c}$ (cert.) & \textsection\ref{sec:qsl} \\
    \bottomrule
  \end{tabular}
\end{table}

The third property is the deep algebraic content of this paper.

\subsection{Existing Responses and Their Limitations}
\label{sec:intro:existing}

Several prior approaches address non-exact framing:

\paragraph{Batz et al.~\cite{batz2019quantitative}.}
Quantitative separation logic (QSL) introduces a $*_\varepsilon$
predicate that allows a bounded additive defect in the separating
conjunction.  Their approach is sound but does not identify the
algebraic origin of the defect---the cocycle class---nor does it provide
a computable bound for general resource monoids.

\paragraph{Aguirre et al.~\cite{aguirre2021general}.}
Graded type theories track effect grades but do not model resource
interference between grades.  The interaction between two composed grades
is not studied.

\paragraph{Probabilistic SL.}
Barthe et al.~\cite{barthe2019probabilistic} handle probability as a
first-class concern but do not classify interference algebraically.

\paragraph{Ghost state and invariants.}
Iris~\cite{jung2018iris} handles subtle invariants via ghost state and
higher-order predicates, but the ghost-state machinery does not give a
direct quantitative bound on frame defects.

None of these gives a \emph{uniform algebraic account} explaining why
frame defects are always compositional (additive under sequential
composition), and none provides a complexity-theoretic analysis of
equivalence modulo bounded defect.

\subsection{This Paper: The Interference Cocycle}
\label{sec:intro:this}

We observe that the resource monoid $(I, \oplus, e)$ of any separation
algebra carries a natural obstruction to additive framing once an
\emph{interference measure} $\rs$ is given.

\begin{definition}[Interference cochain, informal]
  $c(a, b, d) := \rs(a \oplus b, d) - \rs(a, d) - \rs(b, d)$.
\end{definition}

The quantity $c(a, b, d)$ measures how much the combined resource $a
\oplus b$ interferes with context $d$ beyond the sum of individual
interferences.  When $c \equiv 0$, composition is additive and the
standard frame rule holds exactly.  When $c \neq 0$, the classical frame
rule fails, but the defect is exactly $\norm{c}$.

Our key mathematical discovery:
\begin{center}
  \emph{For any monoid $(I, \oplus, e)$ and any $\rs$ additive in its
  second argument, $c$ is always a Hochschild 2-cocycle: $\delta^2 c = 0$.}
\end{center}
This is \emph{not a design choice}: it is forced by monoid associativity
and the right-linearity of $\rs$.  The Hochschild 2-cocycle condition
severely constrains how interference distributes across sequential
compositions, implying in particular that defects accumulate additively
(not multiplicatively) under $;$.  This is the algebraic explanation for
the compositional property (3) above.

\subsection{Summary of Contributions}
\label{sec:intro:contributions}

\begin{enumerate}[label=(\arabic*),leftmargin=2em,itemsep=2pt]
  \item \textbf{Interference cochain and its Hochschild 2-cocycle identity}
    (\S\ref{sec:background}): definitions, the right-linearity axiom, and
    the central identity.  Lean proof: \lean{cocycleClosed} in
    \lean{lean/IdeaTheory/Cocycle.lean}, 9 lines, zero \lean{sorry}.

  \item \textbf{The logic \qSL} (\S\ref{sec:qsl}): syntax, quantitative
    assertions, graded Hoare triples, and the full rule system including
    the cocycle-graded frame rule.

  \item \textbf{Unification of three formalisms}: classical separation
    logic = $c \equiv 0$; graded effects = grade annotation on the resource
    monoid; quantitative type theory = slack bound on the defect.
    Theorem~\ref{thm:triviality} makes this precise.

  \item \textbf{Soundness} (\S\ref{sec:soundness}): the graded frame rule is
    sound with respect to a step-indexed denotational model.
    Theorem~\ref{thm:soundness}.

  \item \textbf{Decidability and PSPACE-completeness} (\S\ref{sec:decidability}):
    finite quotient monoid $M_k$ (Theorem~\ref{thm:finite}), polynomial-space
    decision procedure (Theorem~\ref{thm:decidability}), PSPACE-hardness via
    TQBF reduction (Theorem~\ref{thm:pspace}).

  \item \textbf{Cache-timing case study} (\S\ref{sec:casestudy}):
    a concrete derivation of probabilistic noninterference for AES
    S-box lookup with a computable timing-defect bound
    (Theorem~\ref{thm:ct-ni}).

  \item \textbf{Lean~4 mechanization} (\S\ref{sec:lean}): 600 lines, zero
    \lean{sorry}, covering the cocycle identity, the triple-from-pair
    uniqueness theorem, and two worked instances.
\end{enumerate}

\subsection{Notation Summary}
\label{sec:intro:notation}

Table~\ref{tab:notation} collects the principal symbols used throughout the
paper.  We follow the conventions of O'Hearn et al.~\cite{ohearn2001local}
for separation logic and Hochschild~\cite{hochschild1945cohomology} for
cohomological notation; deviations are noted in the table.

\begin{table}[t]
  \caption{Notation summary.}
  \label{tab:notation}
  \centering\small
  \begin{tabular}{@{}lll@{}}
    \toprule
    Symbol & Meaning & Section \\
    \midrule
    $(M, \oplus, \mathbf{0})$ & commutative resource monoid & \S\ref{sec:background} \\
    $\rs : M \times M \to \mathcal{R}$ & interference/resonance measure & \S\ref{sec:background} \\
    $c(a,b,d)$ & interference 2-cochain & \S\ref{sec:background} \\
    $\delta^2 c$ & Hochschild coboundary & \S\ref{sec:background} \\
    $\HH{2}{M,\mathcal{R}}$ & Hochschild cohomology group & \S\ref{sec:background} \\
    $\norm{c}$ & cocycle norm (sup over $d$) & \S\ref{sec:background} \\
    $\{P\}\;c\;\{Q\}_\delta$ & graded Hoare triple & \S\ref{sec:qsl} \\
    $P \sep Q$ & separating conjunction & \S\ref{sec:qsl} \\
    $P \wand Q$ & separating implication & \S\ref{sec:qsl} \\
    $\fp(c)$ & footprint of command $c$ & \S\ref{sec:qsl} \\
    $\defect$ & frame defect bound & \S\ref{sec:qsl} \\
    $M_k$ & $k$-cocycle quotient monoid & \S\ref{sec:decidability} \\
    $\sim_k$ & $k$-indistinguishability relation & \S\ref{sec:decidability} \\
    $\CTIMP$ & target language (CT-IMP) & \S\ref{sec:qsl} \\
    $\Lambda_{\mathrm{coc}}$ & graded $\lambda$-calculus & \S\ref{sec:qsl} \\
    \bottomrule
  \end{tabular}
\end{table}

\subsection{Separation Algebras}
\label{sec:background:sepalg}

Following Calcagno, O'Hearn, and Yang~\cite{calcagno2007local}, a
\emph{separation algebra} is a structure $(I, \oplus, e)$ where
$\oplus : I \times I \rightharpoonup I$ is a partial binary operation and
$\#(h_1, h_2)$ (read: ``$h_1$ and $h_2$ are compatible'') means
$h_1 \oplus h_2$ is defined.

\begin{definition}[Separation algebra]
\label{def:sepalg}
  $(I, \oplus, e)$ is a \emph{separation algebra} if:
  \begin{itemize}[itemsep=2pt]
    \item (SA-Com) $\#(h_1, h_2) \Rightarrow h_1 \oplus h_2 = h_2 \oplus h_1$.
    \item (SA-Assoc) $\#(h_1, h_2) \land \#(h_1 \oplus h_2, h_3)$
        $\Rightarrow \#(h_2, h_3) \land \#(h_1, h_2 \oplus h_3)$
        $\land (h_1 \oplus h_2) \oplus h_3 = h_1 \oplus (h_2 \oplus h_3)$.
    \item (SA-Unit) $\forall h.\; \#(e, h) \land e \oplus h = h$.
  \end{itemize}
\end{definition}

The canonical instance is the \emph{heap model}: $I = \mathit{Loc}
\rightharpoonup \mathit{Val}$, $\#(h_1, h_2)$ iff $\mathit{dom}(h_1)
\cap \mathit{dom}(h_2) = \emptyset$, and $h_1 \oplus h_2 = h_1 \uplus
h_2$.  Other instances include multiset resources (for counting), formal
power series (for probabilistic resources), and the boolean lattice (for
capability systems).

\begin{remark}[Totalization]
  For the algebraic development of \S\ref{sec:background:hochschild} and
  \S\ref{sec:background:cocycle} we work with the \emph{total} monoid
  $(I, \oplus, e)$ obtained by extending $\oplus$ to a total operation
  via a distinguished ``garbage'' element $\top$ (following
  Dockins et al.~\cite{dockins2009separation}).  All semantic results
  (\S\ref{sec:soundness}) use the partial structure directly.
\end{remark}

\begin{example}[Cache resource monoid]
\label{ex:cache}
  Let $I = \mathcal{P}(\mathit{CacheLine})$ with
  $\oplus = \cup$ and $e = \emptyset$.  This is a total commutative
  monoid that models the set of cache lines occupied by a program
  component.  The compatibility relation is always true (any two
  programs share a cache).  The corresponding interference measure
  $\rs(h, d) = \#(h \cap d) \cdot \Delta$ counts shared lines weighted
  by the miss penalty.
\end{example}

\begin{example}[Probabilistic resource monoid]
\label{ex:prob}
  Let $I = [0,1]$ with $\oplus = \min$ and $e = 1$.  This models
  a \emph{success probability}, where the probability of success of
  two independent components is bounded by the minimum of their
  individual probabilities.  An interference measure
  $\rs(p, q) = p \cdot (1 - q)$ counts the probability that a component
  with probability $p$ fails because a concurrent component has
  failure probability $1 - q$.  The cocycle in this setting measures
  \emph{correlation}: $c(p_1, p_2, q) = \min(p_1, p_2) \cdot (1-q)
  - p_1 (1-q) - p_2 (1-q)$, which is non-zero when $p_1 \neq p_2$
  and $q < 1$.
\end{example}

\begin{example}[Capability resource monoid]
\label{ex:cap}
  Let $I = 2^{\mathit{Cap}}$ (sets of capabilities, e.g., \texttt{read},
  \texttt{write}, \texttt{exec}) with $\oplus = \cup$ and $e = \emptyset$.
  Define $\rs(S, T) = |S \cap T|$ (number of shared capabilities).
  Then $c(A, B, T) = |A \cup B \cap T| - |A \cap T| - |B \cap T|
  = -|A \cap B \cap T|$.  This measures the redundant capability
  contribution: if $A$ and $B$ share a capability $t \in T$, adding
  them does not double the interaction with $T$.
\end{example}

\subsection{Hochschild Cohomology of Monoids}
\label{sec:background:hochschild}

We recall the relevant fragment of Hochschild cohomology~\cite{hochschild1945cohomology}.
Let $(M, \cdot, e)$ be a monoid and $(A, +, 0)$ an abelian group carrying
a \emph{trivial $M$-bimodule structure}: $m \cdot a = a$ and $a \cdot m = a$
for all $m \in M$, $a \in A$.

\begin{definition}[Hochschild cochain complex]
\label{def:hochschild}
  The \emph{normalized Hochschild cochain complex} $C^\bullet(M, A)$ has:
  \begin{itemize}[itemsep=2pt]
    \item $C^n(M, A) = \{ f : M^n \to A \mid f \text{ vanishes when any argument is } e \}$.
    \item Coboundary $\delta^0 : C^0 \to C^1$: $(\delta^0 a)(m) = 0$ (trivial action).
    \item Coboundary $\delta^1 : C^1 \to C^2$:
      $(\delta^1 \varphi)(a, b) = \varphi(b) - \varphi(a \cdot b) + \varphi(a)$.
    \item Coboundary $\delta^2 : C^2 \to C^3$:
      $(\delta^2 f)(a, b, d) = f(b, d) - f(a \cdot b, d) + f(a, b \cdot d) - f(a, b)$.
  \end{itemize}
\end{definition}

The standard identity $\delta^{n+1} \circ \delta^n = 0$ makes
$C^\bullet(M, A)$ a cochain complex.  The cohomology groups
$\HH{n}{M, A} = \ker(\delta^n) / \mathrm{im}(\delta^{n-1})$ are
invariants of the pair $(M, A)$.

For our purposes the critical group is $\HH{2}{M, A}$:

\begin{definition}[2-Cocycles and 2-Coboundaries]
\label{def:2cocycle}
  \begin{itemize}[itemsep=2pt]
    \item $f \in Z^2(M, A)$ (\emph{2-cocycle}) iff $\delta^2 f = 0$, i.e.,
      \[
        f(b, d) - f(a \cdot b, d) + f(a, b \cdot d) - f(a, b) = 0
        \quad \forall a, b, d \in M.
      \]
    \item $f \in B^2(M, A)$ (\emph{2-coboundary}) iff $\exists \varphi \in C^1$
      with $f = \delta^1 \varphi$, i.e., $f(a, b) = \varphi(b) - \varphi(a \cdot b) + \varphi(a)$.
    \item $\HH{2}{M, A} = Z^2 / B^2$.
  \end{itemize}
\end{definition}

\begin{example}[Group extensions]
  For a group $G$ and $G$-module $A$, $\HH{2}{G, A}$ classifies
  central extensions $0 \to A \to \tilde{G} \to G \to 0$.  The
  interference cocycle we introduce is a \emph{monoid} 2-cocycle with
  trivial bimodule structure; its cohomology class classifies
  ``resource monoid extensions by additive defect.''
\end{example}

\begin{remark}[Interpretation in deformation theory]
  In deformation theory, $\HH{2}{A, A}$ classifies first-order
  deformations of an algebra $A$.  Our setting is different: the
  ``algebra'' being deformed is the resource accounting rule (the frame
  rule), and the ``deformation'' is the perturbation of additivity
  captured by $c$.  A trivial deformation ($[c] = 0$) means the frame
  rule holds exactly; a non-trivial deformation means the frame rule
  admits a computable correction $\norm{c}$.
\end{remark}

\begin{remark}[Comparison with associativity cochains]
  The Hochschild 3-cocycle condition $\delta^3 f = 0$ appears in
  MacLane's coherence theorem for monoidal categories.  Our setting uses
  the \emph{2-cocycle} (a 3-variable function $c(a,b,d)$), which arises
  from the right-linearity of $\rs$ and is one degree lower.  This is
  not an accident: the right-linearity of $\rs$ plays the role of the
  ``bimodule condition'' in the Hochschild complex, and it is precisely
  this condition that forces $c$ to be a 2-cocycle rather than a free
  3-cochain.
\end{remark}

\subsection{The Interference Cochain is Always a Cocycle}
\label{sec:background:cocycle}

\begin{definition}[Interference measure]
\label{def:rs}
  An \emph{interference measure} on $(M, \oplus, e)$ valued in an
  abelian group $(R, +, 0)$ is a map $\rs : M \times M \to R$
  satisfying the \emph{right-linearity axiom} (RL):
  \[
    \rs(m,\; k_1 \oplus k_2) = \rs(m, k_1) + \rs(m, k_2)
    \quad \text{for all } m, k_1, k_2 \in M. \tag{RL}
  \]
  Left-linearity $\rs(m_1 \oplus m_2, k) = \rs(m_1, k) + \rs(m_2, k)$
  is intentionally \emph{not} assumed; its failure is measured by $c$.
\end{definition}

\begin{definition}[Interference cochain]
\label{def:cochain}
  Given $(M, \oplus, e, \rs, R)$, the \emph{interference cochain}
  is the map $c : M \times M \times M \to R$ defined by:
  \[
    c(a, b, d) \;:=\; \rs(a \oplus b,\; d) \;-\; \rs(a, d) \;-\; \rs(b, d).
  \]
\end{definition}

Note that $c(a, b, d) = 0$ iff running resources $a$ and $b$ together is
resource-additive for context $d$.  The classical SL disjointness condition
($\#(m, k)$ forces $\rs(m \oplus k, d) = \rs(m, d) + \rs(k, d)$) is exactly
$c = 0$ in the heap-region instance.

\begin{example}[Computing $c$ for cache resource monoid]
\label{ex:cache-c}
  In Example~\ref{ex:cache}, with $I = \mathcal{P}(\mathit{CacheLine})$
  and $\rs(h, d) = |h \cap d| \cdot \Delta$:
  \begin{align*}
    c(h_1, h_2, d) &= |(h_1 \cup h_2) \cap d| \cdot \Delta
                      - |h_1 \cap d| \cdot \Delta - |h_2 \cap d| \cdot \Delta \\
    &= (|h_1 \cap d| + |h_2 \cap d| - |h_1 \cap h_2 \cap d|) \cdot \Delta \\
    &\quad\quad - |h_1 \cap d| \cdot \Delta - |h_2 \cap d| \cdot \Delta \\
    &= -|h_1 \cap h_2 \cap d| \cdot \Delta.
  \end{align*}
  This is always $\leq 0$: cache sharing is a \emph{bonus} (reduces total
  miss count).  $c = 0$ iff $h_1$ and $h_2$ have disjoint intersection
  with context $d$.
\end{example}

\begin{example}[Computing $c$ for probability resource monoid]
\label{ex:prob-c}
  In Example~\ref{ex:prob}, with $I = [0,1]$, $\oplus = \min$, and
  $\rs(p, q) = p \cdot (1-q)$:
  \begin{align*}
    c(p_1, p_2, q) &= \rs(\min(p_1, p_2), q) - \rs(p_1, q) - \rs(p_2, q) \\
    &= \min(p_1, p_2)(1-q) - p_1(1-q) - p_2(1-q) \\
    &= (\min(p_1, p_2) - p_1 - p_2)(1-q).
  \end{align*}
  This is $\leq 0$ (since $\min(p_1, p_2) \leq p_1$ and $\leq p_2$, so the
  sum $\min - p_1 - p_2 \leq -\max(p_1, p_2) \leq 0$), capturing the
  probability mass ``lost'' by the min operation.
\end{example}

The main algebraic theorem of the paper:

\begin{theorem}[Cocycle identity; \leancite{cocycleClosed}{Cocycle.lean}, line 74]
\label{thm:cocycle}
  For any monoid $(M, \oplus, e)$ and interference measure $\rs$
  satisfying \emph{(RL)}, the interference cochain $c$ is a Hochschild
  2-cocycle:
  \[
    (\delta^2 c)(a, b, d) \;=\; 0 \quad \text{for all } a, b, d \in M.
  \]
  Equivalently:
  \[
    c(b, d) - c(a \oplus b, d) + c(a, b \oplus d) - c(a, b) = 0.
  \]
\end{theorem}

\begin{proof}
  Expand each of the four $c$-terms into $\rs$-terms.  Specifically:
  \begin{align*}
    c(b, d) &= \rs(b \oplus d, \cdot) - \rs(b, \cdot) - \rs(d, \cdot), \\
    c(a \oplus b, d) &= \rs((a \oplus b) \oplus d, \cdot) - \rs(a \oplus b, \cdot) - \rs(d, \cdot), \\
    c(a, b \oplus d) &= \rs(a \oplus (b \oplus d), \cdot) - \rs(a, \cdot) - \rs(b \oplus d, \cdot), \\
    c(a, b) &= \rs(a \oplus b, \cdot) - \rs(a, \cdot) - \rs(b, \cdot).
  \end{align*}
  (Here all $\rs$-terms are evaluated with third argument from context;
  we suppress it for readability.)
  After cancellation:
  \begin{itemize}[itemsep=1pt]
    \item $\rs((a \oplus b) \oplus d, \cdot)$ and $\rs(a \oplus (b \oplus d), \cdot)$
      are equal by \emph{associativity of $\oplus$}, and appear with
      opposite signs $-$ and $+$.
    \item $\rs(b \oplus d, \cdot)$ appears with signs $+$ (from $c(b,d)$)
      and $-$ (from $c(a, b \oplus d)$); cancels.
    \item $\rs(a \oplus b, \cdot)$ appears with signs $+$ (from $c(a \oplus b, d)$,
      note the leading $-$ flips) and $-$ (from $c(a, b)$); cancels.
    \item $\rs(a, \cdot)$, $\rs(b, \cdot)$, $\rs(d, \cdot)$ each appear
      exactly twice with opposite signs and cancel.
  \end{itemize}
  All eight $\rs$-terms cancel, giving 0.  The Lean proof formalizes this
  by \lean{rw [add\_assoc a b d]}, three applications of
  \lean{P.add\_right} (for the (RL) axiom in additive notation), and
  \lean{abel}.
\end{proof}

\begin{corollary}[Cohomological obstruction]
\label{cor:coh}
  The cohomology class $[c] \in \HH{2}{M, R}$ is an invariant of
  $(M, \oplus, \rs)$.  $[c] = 0$ iff $\rs(\cdot, d)$ is a monoid
  homomorphism $(M, \oplus) \to (R, +)$ for every $d$.
\end{corollary}

\begin{corollary}[Frame rule iff cocycle triviality]
\label{cor:frame-iff-cocycle}
  The standard frame rule holds for resource measure $\rs$ (meaning:
  adding any framing resource $R$ to any command's pre/postcondition
  introduces zero defect) if and only if $[c] = 0$.
\end{corollary}

\subsection{The Cocycle Norm and Defect Bound}
\label{sec:background:norm}

The key invariant for the program logic is the \emph{cocycle norm}:

\begin{definition}[Cocycle norm]
\label{def:cnorm}
  $\norm{c}_{f, r} = \sup\{ |c(a, b, d)| \mid a \in f,\; b \in r,\; d \in I \}$
  where $f \subseteq I$ is the command footprint and $r \subseteq I$ is
  the frame support.
\end{definition}

\begin{lemma}[Norm monotonicity]
\label{lem:norm-mono}
  If $f \subseteq f'$ and $r \subseteq r'$, then
  $\norm{c}_{f, r} \leq \norm{c}_{f', r'}$.
\end{lemma}

\begin{lemma}[Norm under sequential composition]
\label{lem:norm-seq}
  For sequential composition $c_1; c_2$ with $\fp(c_1) = f_1$ and
  $\fp(c_2) = f_2$, the defect of the composed command framed against $r$ is:
  \[
    \norm{c}_{f_1 \cup f_2, r} \leq \norm{c}_{f_1, r} + \norm{c}_{f_2, r}.
  \]
\end{lemma}
\begin{proof}
  For any $a \in f_1 \cup f_2$, either $a \in f_1$ or $a \in f_2$
  (or both).  If $a \in f_1$: $|c(a, b, d)| \leq \norm{c}_{f_1, r}$ for
  any $b \in r$, $d \in I$.  Similarly for $f_2$.  Taking the sup:
  $\norm{c}_{f_1 \cup f_2, r} \leq \max(\norm{c}_{f_1, r}, \norm{c}_{f_2, r})
  \leq \norm{c}_{f_1, r} + \norm{c}_{f_2, r}$.
\end{proof}

\begin{remark}[When does the cocycle norm vanish?]
  $\norm{c}_{f, r} = 0$ iff for all $a \in f$, $b \in r$, $d \in I$:
  $\rs(a \oplus b, d) = \rs(a, d) + \rs(b, d)$.  This is the condition
  that $a$ and $b$ are \emph{resource-additive}: their combined effect is
  the sum of individual effects.  For the heap model, this holds whenever
  $\mathrm{dom}(a) \cap \mathrm{dom}(b) = \emptyset$, i.e., when the
  command's footprint and the frame are disjoint.
\end{remark}

\begin{definition}[Pairing rank]
\label{def:rank}
  The \emph{pairing rank} $r$ of $(M, \rs)$ is the rank of the linear map
  $a \mapsto (\rs(a, d_1), \ldots, \rs(a, d_m))$ from $M$ to $\mathbb{Z}^m$
  (where $d_1, \ldots, d_m$ are the elements of $M$).  The rank measures
  how many independent interference ``dimensions'' $\rs$ has.
\end{definition}

The pairing rank controls the complexity of the quotient monoid $M_k$
(Theorem~\ref{thm:finite}) and hence the complexity of the decision
procedure (Theorem~\ref{thm:decidability}).  For the cache-timing monoid
$(\mathcal{P}(\mathit{CacheLine}), \cup)$, the pairing rank equals the
number of cache lines (since each line contributes an independent dimension
to the interference profile).

\begin{example}[Non-trivial cocycle on $\mathbb{Z}^2$]
\label{ex:quadratic}
  Define $\rs : (\mathbb{Z}^2, +, \mathbf{0}) \times (\mathbb{Z}^2, +, \mathbf{0})
  \to \mathbb{Z}$ by
  $\rs((a_1, a_2), (b_1, b_2)) = a_1^2 b_1 + a_2 b_2$.
  This is additive in the second argument (right-linear) but \emph{not}
  additive in the first (since $(a_1 + a_1')^2 \neq a_1^2 + (a_1')^2$ in
  general).  The emergence cochain is:
  \begin{align*}
    c((1,0),(1,0),(1,0)) &= \rs((2,0),(1,0)) - \rs((1,0),(1,0)) - \rs((1,0),(1,0)) \\
    &= 4 \cdot 1 - 1 \cdot 1 - 1 \cdot 1 = 2.
  \end{align*}
  So $c \neq 0$, but Theorem~\ref{thm:cocycle} guarantees $\delta^2 c = 0$.
  This is verified by \lean{native\_decide} in the Lean artifact.
\end{example}

\begin{example}[IO read/write effects]
\label{ex:io}
  Let $M = \{\varepsilon, \mathit{rd}, \mathit{wr}, \mathit{rdwr}\}$
  with the quotient monoid structure $\mathit{rd} \oplus \mathit{wr}
  = \mathit{wr} \oplus \mathit{rd} = \mathit{rdwr}$ (free with some
  equations).  Define $\rs(m, k)$ = number of bytes read or written to
  file $k$ by effect $m$.  Then:
  \[
    c(\mathit{rd}, \mathit{wr}, k) = \rs(\mathit{rdwr}, k) - \rs(\mathit{rd}, k) - \rs(\mathit{wr}, k).
  \]
  This is the \emph{read-write interference}: bytes whose write target
  was chosen because of the read.  The cocycle identity constrains how
  this interference propagates through longer compositions: the
  interference of $(\mathit{rd} \oplus \mathit{wr}) \oplus m_3$ with
  $k$ is determined by the pairwise interferences alone.
\end{example}

\begin{example}[Security lattice]
\label{ex:sec}
  Let $M = (\mathcal{L}, \vee, \bot)$ be a security lattice (join
  semilattice).  Define $\rs(\ell_1, \ell_2) = [\ell_1 \geq \ell_2]$
  (Boolean: does $\ell_1$ flow to $\ell_2$?).  For a non-trivial lattice
  with $H > L$, $c(H, L, M) = \rs(H \vee L, M) - \rs(H, M) - \rs(L, M)
  = \rs(H, M) - \rs(H, M) - \rs(L, M) = -\rs(L, M)$.
  This is non-zero when $L$ flows to $M$ (i.e., $M \geq L$), capturing
  the covert channel created by combining a high-security effect with a
  low-security one.
\end{example}

Table~\ref{tab:instances} lists concrete interference measure instances.

\begin{table}[t]
  \caption{Instances of the interference measure.}
  \label{tab:instances}
  \centering\small
  \begin{tabular}{@{}p{2cm}p{2cm}p{2.5cm}p{1.5cm}@{}}
    \toprule
    System & Monoid $M$ & $\rs(m, k)$ & $c$ \\
    \midrule
    Classical SL & Heap regions & $[m \cap k \neq \emptyset]$ & $0$ \\
    Linear types & $(\mathbb{N}, +, 0)$ & $m \cdot k$ & $0$ \\
    Cache timing & $\mathcal{P}(\text{lines})$ & $|m \cap k| \cdot \Delta$ & $\leq 0$ \\
    Info flow & Security lattice & $\mathrm{flow}(m \to k)$ & $\neq 0$ \\
    Prob. sched. & Distributions & $\mathrm{Cov}(m, k)$ & $\neq 0$ \\
    IO effects & $\{\varepsilon,\mathit{rd},\mathit{wr},\mathit{rdwr}\}$ & byte overlap & $\neq 0$ \\
    $\mathbb{Z}^2$ quadratic & $(\mathbb{Z}^2, +)$ & $a_1^2 b_1 + a_2 b_2$ & $\neq 0$ \\
    \bottomrule
  \end{tabular}
\end{table}

% ==================================================================
\section{The \qSL Calculus}
\label{sec:qsl}

\subsection{The Target Language: CT-IMP}
\label{sec:qsl:language}

We define \CTIMP (Cache-Timing IMP), a concrete target language for the
case study (\S\ref{sec:casestudy}).  The logic \qSL is parametric in the
resource monoid and can be instantiated for other languages.

\paragraph{Syntax.}
\begin{align*}
  \text{Expr}\;e &::= n \mid x \mid e_1 \mathbin{\mathit{op}} e_2 \\
  \text{Cmd}\;c  &::= \mathbf{skip} \mid x \mathbin{:=} e \mid c_1 ; c_2
                      \mid \mathbf{if}\; e\; \mathbf{then}\; c_1\; \mathbf{else}\; c_2 \\
                 &\phantom{::=}\mid \mathbf{while}\; e\; \mathbf{do}\; c
                      \mid \mathbf{load}\; x \mathbin{\leftarrow} [e]
                      \mid \mathbf{store}\; [e] \mathbin{\leftarrow} e'
\end{align*}

\paragraph{State.}
A \CTIMP \emph{state} is a triple $(\sigma, h, \kappa)$ where:
$\sigma : \mathit{Var} \to \mathit{Val}$ is the variable store,
$h : \mathit{Loc} \rightharpoonup \mathit{Val}$ is the heap,
$\kappa : \mathit{CacheLine} \to \{\mathit{hot}, \mathit{cold}\}$ is the
cache state.

\paragraph{Small-step semantics.}
The small-step relation $\langle c, \sigma, h, \kappa \rangle
\xrightarrow{t} \langle c', \sigma', h', \kappa' \rangle$ is annotated
with timing cost $t \in \mathbb{N}$.  The cache-timing rules are:
\begin{itemize}[itemsep=2pt]
  \item \textbf{Cache-hit load} (address $a = \llbracket e \rrbracket_\sigma$ maps to hot line
    $\ell$):
    $\langle \mathbf{load}\; x \leftarrow [e], \sigma, h, \kappa \rangle
    \xrightarrow{t_{\mathrm{hit}}}
    \langle \mathbf{skip}, \sigma[x \mapsto h(a)], h, \kappa \rangle$.
  \item \textbf{Cache-miss load} (line $\ell$ cold, LRU evicts line $\ell'$):
    $\langle \mathbf{load}\; x \leftarrow [e], \sigma, h, \kappa \rangle
    \xrightarrow{t_{\mathrm{miss}}}
    \langle \mathbf{skip}, \sigma[x \mapsto h(a)], h, \kappa[\ell \mapsto \mathit{hot}, \ell' \mapsto \mathit{cold}] \rangle$.
  \item All other commands have cost $t = 1$.
\end{itemize}

Let $\Delta := t_{\mathrm{miss}} - t_{\mathrm{hit}}$ be the \emph{miss penalty}.

\paragraph{Operational semantics for loops.}
The semantics of the \textbf{while} command uses the standard loop-unrolling
rule parameterized by a step index $n$, and the timing annotation accumulates:
\begin{align*}
  &\langle \mathbf{while}\; e\; \mathbf{do}\; c, \sigma, h, \kappa \rangle
  \;\xrightarrow{1}\; \\
  &\quad\begin{cases}
    \langle \mathbf{skip}, \sigma, h, \kappa \rangle
    & \text{if } \llbracket e \rrbracket_\sigma = 0, \\
    \langle c ; \mathbf{while}\; e\; \mathbf{do}\; c, \sigma, h, \kappa \rangle
    & \text{otherwise.}
  \end{cases}
\end{align*}
Timing cost 1 is charged for the condition evaluation regardless of the
branch taken.  The total timing cost of $n$ loop iterations is therefore
$n \cdot T_c + n$ where $T_c$ is the per-iteration cost of $c$.

\paragraph{Footprint.}
The \emph{static footprint} $\fp(c) \in \mathcal{P}(\mathit{CacheLine})$
is defined inductively:
\begin{align*}
  \fp(\mathbf{skip}) &= \emptyset, &
  \fp(x := e) &= \emptyset, \\
  \fp(c_1; c_2) &= \fp(c_1) \cup \fp(c_2), \\
  \fp(\mathbf{load}\; x \leftarrow [e]) &= \{\mathrm{line}(e)\}, \\
  \fp(\mathbf{store}\; [e] \leftarrow e') &= \{\mathrm{line}(e)\}, \\
  \fp(\mathbf{if}\; e\; \mathbf{then}\; c_1\; \mathbf{else}\; c_2) &= \fp(c_1) \cup \fp(c_2), \\
  \fp(\mathbf{while}\; e\; \mathbf{do}\; c) &= \fp(c),
\end{align*}
where $\mathrm{line}(e)$ maps an address expression to its cache line.
The footprint is a \emph{conservative over-approximation}: $c$ may access
fewer cache lines at runtime (e.g., loop may not execute), but never more.
For the timing defect bound, using the static footprint is sound and yields
the tightest certifiable bound derivable from syntactic analysis alone.

\begin{lemma}[Footprint conservatism]
\label{lem:fp-conserv}
  For all states $(\sigma, h, \kappa)$ and all $c$:
  $\mathrm{accessed}(c, \sigma, h) \subseteq \fp(c)$,
  where $\mathrm{accessed}(c, \sigma, h)$ is the set of cache lines
  accessed during any terminating execution of $c$.
\end{lemma}

\begin{proof}
  By structural induction on $c$.  The key case is \textbf{load}: the
  accessed line is $\mathrm{line}(\llbracket e \rrbracket_\sigma)$, which
  is a runtime value; but $\mathrm{line}(e) \in \fp(\mathbf{load}\; x \leftarrow [e])$
  by definition, and by assumption $\mathrm{line}(e)$ is the static
  approximation (i.e., for array-type programs, the base address is
  statically known).  The other cases follow trivially.
\end{proof}

\paragraph{Footprint.}
$\fp(c, \sigma, h)$ is the set of heap locations accessed by $c$ in state
$(\sigma, h)$.  We write $\fp(c)$ when the state is clear from context.

\subsection{Quantitative Assertions}
\label{sec:qsl:assertions}

Fix a \emph{normed semiring} $(\mathcal{R}, +, \cdot, \norm{\cdot}, 0, 1)$
with $\norm{\cdot}$ satisfying $\norm{a + b} \leq \norm{a} + \norm{b}$
and $\norm{a \cdot b} \leq \norm{a} \cdot \norm{b}$.

\begin{definition}[Quantitative assertions]
\label{def:qassert}
  \emph{Quantitative assertions} are functions $P : I \to \mathcal{R}$.
  \begin{align*}
    \emp(h) &\;:=\; [h = e], \\
    \lceil q \rceil(h) &\;:=\; q \quad \text{(constant)}, \\
    (P \sep Q)(h) &\;:=\; \sup\{ P(h_1) \cdot Q(h_2) \mid \#(h_1, h_2),\; h_1 \oplus h_2 = h \}, \\
    (P \wand Q)(h) &\;:=\; \inf\{ Q(h \oplus h') / P(h') \mid \#(h, h'),\; P(h') > 0 \}.
  \end{align*}
\end{definition}

The separating conjunction $P \sep Q$ assigns to $h$ the maximum
``combined value'' over all compatible splits.  For Boolean predicates
(where $\mathcal{R} = \{0,1\}$) this reduces to the classical
$\exists h_1, h_2.\; h_1 \oplus h_2 = h \land P(h_1) \land Q(h_2)$.

\subsection{Graded Hoare Triples}
\label{sec:qsl:triples}

\begin{definition}[Graded Hoare triple]
\label{def:triple}
  $\vdash \{P\}\; c\; \{Q\}_\defect$ (\emph{graded Hoare triple}) with
  \emph{defect bound} $\defect \in \mathcal{R}_{\geq 0}$ is well-formed
  when the following \emph{semantic validity} holds (see
  \S\ref{sec:semantics} for the formal definition):
  \[
    \forall h.\; P(h) \;\leq\; Q(\llbracket c \rrbracket(h)) + \defect.
  \]
  The zero-defect case $\defect = 0$ coincides with standard quantitative
  Hoare triples~\cite{aguirre2021general}.
\end{definition}

\subsection{Proof Rules}
\label{sec:qsl:rules}

The rules of \qSL are given in Figure~\ref{fig:rules}.

\begin{figure}[t]
  \hrule\smallskip
  \paragraph{Structural rules.}
  \[
    \inferrule{ }{\vdash \{\,P\,\}\;\mathbf{skip}\;\{\,P\,\}_0}
    \quad\text{[Skip]}
  \]
  \[
    \inferrule{P' \leq P \quad \vdash \{P\}\;c\;\{Q\}_\defect \quad Q \leq Q' + \defect'}
              {\vdash \{P'\}\;c\;\{Q'\}_{\defect + \defect'}}
    \quad\text{[Conseq]}
  \]
  \[
    \inferrule{\vdash \{P\}\;c_1\;\{R\}_{\defect_1} \quad \vdash \{R\}\;c_2\;\{Q\}_{\defect_2}}
              {\vdash \{P\}\;c_1;c_2\;\{Q\}_{\defect_1+\defect_2}}
    \quad\text{[Seq]}
  \]

  \paragraph{Command rules.}
  \[
    \inferrule{ }{\vdash \{P[v \mapsto e]\}\;x \mathbin{:=} e\;\{P\}_0}
    \quad\text{[Assign]}
  \]
  \[
    \inferrule{\vdash \{P \land b\}\;c_1\;\{Q\}_{\defect_1} \quad
               \vdash \{P \land \neg b\}\;c_2\;\{Q\}_{\defect_2}}
              {\vdash \{P\}\;\mathbf{if}\;b\;\mathbf{then}\;c_1\;\mathbf{else}\;c_2\;\{Q\}_{\max(\defect_1,\defect_2)}}
    \quad\text{[If]}
  \]
  \[
    \inferrule{\vdash \{I \land b\}\;c\;\{I\}_{\defect}}
              {\vdash \{I\}\;\mathbf{while}^n\;b\;\mathbf{do}\;c\;\{I \land \neg b\}_{n \cdot \defect}}
    \quad\text{[While]}
  \]

  \paragraph{The cocycle-graded frame rule.}
  \[
    \inferrule{\vdash \{P\}\;c\;\{Q\}_\defect}
              {\vdash \{P \sep R\}\;c\;\{Q \sep R\}_{\defect + \defect_{\mathrm{frame}}}}
    \quad\text{[GrFrame]}
  \]
  where
  $\defect_{\mathrm{frame}} = \sup\,\{|c(a, b, d)| \mid a \in \fp(c),\;
  b \in \mathrm{supp}(R),\; d \in I\}$.\smallskip

  \paragraph{Parallel composition (concurrent extension).}
  \[
    \inferrule{\vdash \{P_1\}\;c_1\;\{Q_1\}_{\defect_1} \quad \vdash \{P_2\}\;c_2\;\{Q_2\}_{\defect_2}
               \quad \defect_{12} \geq \mathrm{cross}(c_1, c_2)}
              {\vdash \{P_1 \sep P_2\}\;c_1 \| c_2\;\{Q_1 \sep Q_2\}_{\defect_1 + \defect_2 + \defect_{12}}}
    \quad\text{[Par]}
  \]
  where $\mathrm{cross}(c_1, c_2) = \sup_{h,h'} |\rs(\fp(c_1)(h), \fp(c_2)(h'))|$.\medskip
  \hrule
  \caption{Proof rules of \qSL.  The defect $\defect_{\mathrm{frame}}$ in
    \textbf{[GrFrame]} is computed from the interference cochain $c$.}
  \label{fig:rules}
\end{figure}

\paragraph{Discussion of [GrFrame].}
The frame defect $\defect_{\mathrm{frame}}$ depends on:
\begin{enumerate}[label=(\roman*),itemsep=2pt]
  \item The footprint $\fp(c)$ of the command,
  \item The support $\mathrm{supp}(R)$ of the frame assertion, and
  \item The full resource monoid $I$ via the context argument $d$.
\end{enumerate}
For the cache-timing instance (Example~\ref{ex:cache}),
$\defect_{\mathrm{frame}} = |\fp(c) \cap \mathrm{supp}(R)| \cdot \Delta$,
which is the number of shared cache lines times the miss penalty---exactly
the expected timing leak.  When $\fp(c) \cap \mathrm{supp}(R) = \emptyset$
(disjoint footprints), $\defect_{\mathrm{frame}} = 0$ and
\textbf{[GrFrame]} becomes the classical frame rule.

\begin{example}[Frame defect for three-command sequence]
\label{ex:frame3}
  Consider three commands $c_1; c_2; c_3$ with footprints $f_1, f_2, f_3$
  and a frame $R$ with support $r$.  By three applications of \textbf{[Seq]}
  and \textbf{[GrFrame]}:
  \[
    \vdash \{P \sep R\}\; c_1;c_2;c_3\; \{Q \sep R\}_{\defect_1 + \defect_2 + \defect_3 + \defect_R}
  \]
  where:
  \begin{align*}
    \defect_i &= \sup\{|c(a, b, d)| \mid a \in f_i,\; b \in \mathit{supp}(\text{intermediate}),\; d \in I\}, \\
    \defect_R &= \sum_{i=1}^3 \sup\{|c(a, b, d)| \mid a \in f_i,\; b \in r,\; d \in I\}.
  \end{align*}
  The frame defect $\defect_R$ is the sum of individual frame defects, not
  the product or some larger quantity.  This is the direct consequence of
  $\delta^2 c = 0$: the interaction of each $f_i$ with $r$ is independent
  (no multiplicative blow-up from chaining).
\end{example}

\begin{example}[Constant-time program (zero defect)]
\label{ex:const-time}
  A program $c$ is \emph{timing-constant} with respect to a separation
  algebra if $c(a, b, d) = 0$ for all $a \in \fp(c)$ and $b \in r$ (any
  frame resource $r$).  In this case \textbf{[GrFrame]} with $\defect_{\mathrm{frame}}=0$
  gives:
  \[
    \vdash \{P \sep R\}\; c\; \{Q \sep R\}_0
  \]
  recovering the exact classical frame rule.  The decision procedure of
  \S\ref{sec:decidability} checks whether a program is timing-constant in
  time $O(n + m^4 \log B)$, providing an automated certificate.
\end{example}

\begin{example}[Loop with frame accumulation]
\label{ex:loop-frame}
  Let $c_{\mathrm{body}}$ be a loop body with $\fp(c_{\mathrm{body}}) = f$
  and $\defect_{\mathrm{frame}} = k$ per iteration.  For
  $\mathbf{while}^n\; b\; \mathbf{do}\; c_{\mathrm{body}}$, the
  \textbf{[While]} rule gives defect $n \cdot k$ at most $n$ iterations.
  Combined with \textbf{[GrFrame]} for a frame $R$ with $\mathrm{supp}(R) \cap f \neq \emptyset$:
  \[
    \vdash \{I \sep R\}\; \mathbf{while}^n\; b\; \mathbf{do}\; c_{\mathrm{body}}\; \{I' \sep R\}_{n(k + \defect_R)}
  \]
  The total defect $n(k + \defect_R)$ grows linearly with the iteration
  count.  For a fixed iteration bound $n$ (e.g., a loop over 256 S-box
  entries), this gives a computable finite defect for the entire loop.
\end{example}

\subsection{Correctness: the Cocycle Frame Rule Implies Classical SL}
\label{sec:qsl:classical}

\begin{corollary}[Classical separation logic as a special case]
\label{cor:sl}
  If $M$ is a separation algebra with $c \equiv 0$ (i.e., the frame rule
  holds exactly), then every derivation in \qSL with all defects $= 0$ is
  also a valid derivation in classical separation logic (Reynolds'
  \textbf{Frame} rule).
\end{corollary}
\begin{proof}
  Setting $\defect_{\mathrm{frame}} = 0$ in \textbf{[GrFrame]} recovers
  the classical rule.  The other rules reduce to their classical
  counterparts when all defect bounds are 0.
\end{proof}

\begin{corollary}[Graded effect systems as a special case]
\label{cor:graded}
  If $\defect_{\mathrm{frame}} = 0$ for all command/frame pairs (but
  the effect grade annotation $m \in M$ is tracked), then \qSL reduces
  to a graded effect system in the sense of Katsumata~\cite{katsumata2014parametric}:
  $\vdash \{P\}\; c\; \{Q\}_m$ with grade $m$ replacing the defect.
\end{corollary}
\begin{proof}
  The grade $m$ records the resource usage; the defect is 0 because the
  frame assertions are disjoint from the command's footprint in a graded
  system.  The [GrFrame] rule specializes to graded frame rule of Katsumata.
\end{proof}

\subsection{The Graded $\lambda$-Calculus $\Lambda_{\mathrm{coc}}$}
\label{sec:qsl:lambda}

To further motivate the logic, we present $\Lambda_{\mathrm{coc}}$, a
simply-typed graded $\lambda$-calculus whose type system tracks both the
resource grade and the cocycle defect.  This calculus is a surface syntax
for \qSL and illustrates how graded types and the quantitative frame rule
arise from a common source.

\paragraph{Types and terms.}
\begin{align*}
  \text{Types}\; A, B &::= \mathbf{base} \mid A \xrightarrow{m} B \mid \Box^m A \\
  \text{Terms}\; e &::= x \mid \lambda x.\, e \mid e_1\; e_2 \mid
    \mathbf{return}\; e \mid \mathbf{bind}\; x = e_1\; \mathbf{in}\; e_2
\end{align*}
The annotation $m \in M$ on arrows records the \emph{effect grade} of the
function body; $\Box^m A$ is the type of computations with grade $m$ that
produce an $A$.

\paragraph{Graded typing judgements.}
A \emph{graded typing judgement} has the form
$\Gamma \vdash e : A \mid m$ where $\Gamma = x_1 :_{m_1} A_1, \ldots,
x_n :_{m_n} A_n$ is a resource-graded context.

Selected rules:
\[
  \inferrule{ }{\Gamma, x :_m A \vdash x : A \mid m}
  \quad\text{[Var]}
\]
\[
  \inferrule{\Gamma \vdash e_1 : A \xrightarrow{m_2} B \mid m_1 \quad
             \Gamma \vdash e_2 : A \mid m_3}
            {\Gamma \vdash e_1\; e_2 : B \mid m_1 \oplus m_2 \oplus m_3}
  \quad\text{[App]}
\]
\[
  \inferrule{\Gamma \vdash e : A \mid m}
            {\Gamma, x :_k B \vdash e : A \mid m}
  \quad \text{with slack } \|c(m, k)\|
  \quad\text{[Weak]}
\]
The [Weak] rule corresponds to the frame rule: adding a resource $k$ to
context introduces a defect $\|c(m, k)\| = \sup_d |c(m, k, d)|$.  When
$c \equiv 0$ (the interference measure is left-linear), [Weak] has zero
defect, recovering standard weakening.

\paragraph{Type safety.}
The progress and preservation theorems for $\Lambda_{\mathrm{coc}}$ hold
in the standard form, augmented by a \emph{defect monotonicity} property:
if $\Gamma \vdash e : A \mid m$ and $e \to e'$, then
$\Gamma \vdash e' : A \mid m'$ with $\|c(m', k)\| \leq \|c(m, k)\|$
for all $k$ (reductions do not increase defects).

\paragraph{Connection to \qSL.}
The translation from $\Lambda_{\mathrm{coc}}$ to \qSL maps:
\begin{itemize}[itemsep=2pt]
  \item Each graded judgement $\Gamma \vdash e : A \mid m$ to the triple
    $\vdash \{P_\Gamma\}\; \llbracket e \rrbracket\; \{P_A\}_{\defect(m)}$
    where $P_\Gamma$ encodes the resource context and
    $\defect(m) = \sup_{k \in \mathrm{context}} \|c(m, k)\|$.
  \item The [Weak] rule to \textbf{[GrFrame]}.
  \item The [App] rule to \textbf{[Seq]}.
\end{itemize}
This translation is \emph{coherent}: the defect of a composed term
$e_1\; e_2$ equals the sum of the defects of $e_1$ and $e_2$, by the
cocycle identity (Theorem~\ref{thm:cocycle}).

\begin{example}[Type-checking an AES round in $\Lambda_{\mathrm{coc}}$]
\label{ex:lambda-aes}
  Let the resource grades be:
  $m_{\mathrm{idx}} = \{\}$ (assignment, no heap),
  $m_{\mathrm{load}} = \mathrm{sbox}$ (loads from S-box),
  $m_{\mathrm{out}} = \mathrm{out}$ (writes to output).

  The typing derivation for \texttt{aes\_round} in $\Lambda_{\mathrm{coc}}$:
  \[
    x :_{\{\}} \mathsf{Int}, \;\mathit{sbox} :_{\mathrm{sbox}} \mathsf{Arr},\;
    \mathit{out} :_{\mathrm{out}} \mathsf{Ptr}
    \;\vdash\; \texttt{aes\_round} : \mathsf{Int} \mid m_{\mathrm{idx}} \oplus m_{\mathrm{load}} \oplus m_{\mathrm{out}}
  \]
  The combined grade is $m = \mathrm{sbox} \oplus \mathrm{out}$.
  The defect with respect to a frame assertion for \texttt{key} and
  \texttt{data} is:
  \[
    \defect(m) = \sup_{k \in \{\mathit{key}, \mathit{data}\}} \|c(m, k)\|
    = \|c(\mathrm{sbox}, \mathrm{key})\| + \|c(\mathrm{sbox}, \mathrm{data})\|.
  \]
  Since \texttt{key} and \texttt{data} are heap-disjoint from the S-box,
  $c(\mathrm{sbox}, \mathrm{key}) = 0$ and $c(\mathrm{sbox}, \mathrm{data}) = 0$.
  However, in the cache model (where all memory shares cache lines),
  $\|c(\mathrm{sbox}, \mathrm{key})\| = L \cdot \Delta$ (the S-box footprint
  shares up to $L$ cache lines with \texttt{key}).
  The $\Lambda_{\mathrm{coc}}$ typing records this defect:
  $\texttt{aes\_round} : \mathsf{Int} \mid m$ with $\defect(m) = L\Delta$.
\end{example}

\paragraph{Normalization and confluence.}
The $\beta$-reduction rules of $\Lambda_{\mathrm{coc}}$ preserve defects:
if $e \to_\beta e'$, then $\defect(e') \leq \defect(e)$.  This is because
reduction does not change the set of resource grades used; it only
substitutes values.  Confluence holds by the Church--Rosser property
of the underlying simply-typed lambda calculus, since the grade annotations
do not affect reduction.

\subsection{Cohomological Classification}
\label{sec:qsl:cohom}

\begin{theorem}[Frame triviality iff cocycle class]
\label{thm:triviality}
  Let $(I, \oplus, e, \rs)$ be a separation algebra with interference
  measure.  The following are equivalent:
  \begin{enumerate}[label=(\roman*),itemsep=2pt]
    \item The standard zero-defect frame rule
      $\inferrule{\{P\}\;c\;\{Q\}}{\{P \sep R\}\;c\;\{Q \sep R\}_0}$
      is sound for all $c$ and $R$.
    \item For every $d \in I$, the map $\rs(\cdot, d) : (I, \oplus) \to
      (\mathcal{R}, +)$ is a monoid homomorphism.
    \item $[c] = 0 \in \HH{2}{I, \mathcal{R}}$.
  \end{enumerate}
\end{theorem}

\begin{proof}
  (i)$\Rightarrow$(ii): If the frame rule holds with zero defect for
  all commands, take $c$ to be an abstract command with footprint $a \oplus b$
  and frame assertion supported at $d$.  Zero defect means $\rs(a \oplus b, d)
  = \rs(a, d) + \rs(b, d)$, i.e., $\rs(\cdot, d)$ is a homomorphism.

  (ii)$\Rightarrow$(iii): If $\rs(\cdot, d)$ is a homomorphism for all $d$
  then $c(a, b, d) = 0$ for all $a, b, d$, so $c$ is the zero cochain,
  and $[c] = [0] = 0$.

  (iii)$\Rightarrow$(i): $[c] = 0$ means $c = \delta^1 \varphi$ for some
  $\varphi \in C^1$.  Compute: $c(a, b, d) = \varphi(b) \cdot d -
  \varphi(a \cdot b) \cdot d + \varphi(a) \cdot d$.  But $c(a,b,d)$ also
  equals $\rs(a \oplus b, d) - \rs(a,d) - \rs(b,d)$ by definition.  Since
  $d$ is arbitrary, we can conclude $c \equiv 0$ by choosing $d$ to
  separate values.  Hence $\defect_{\mathrm{frame}} = 0$ for all $c$ and $R$.
\end{proof}

% ==================================================================
\section{Denotational and Operational Semantics}
\label{sec:semantics}

\subsection{Step-Indexed Quantitative Model}
\label{sec:semantics:model}

We use a step-indexed semantics to handle non-terminating programs.

\begin{definition}[Step-indexed predicate]
\label{def:si-pred}
  A \emph{step-indexed quantitative predicate} $(P_n)_{n \in \mathbb{N}}$
  is a family $P_n : I \to \mathcal{R}_{\geq 0}$ satisfying
  \emph{downward monotonicity}: $n' \leq n \Rightarrow P_{n'}(h) \leq P_n(h)$.
  The limit $P = \lim_n P_n$ exists in the complete lattice
  $\mathcal{R}_{\geq 0}^I$ since $\mathcal{R}_{\geq 0}$ is complete.
\end{definition}

\begin{definition}[Denotational semantics]
\label{def:denotation}
  The $n$-step denotation
  $\llbracket c \rrbracket_n : I \times \Sigma \times K \to
  (I \times \Sigma \times K \times \mathbb{N}) \cup \{\bot\}$
  is defined by structural induction on $c$, where $\Sigma$ is the
  store space and $K$ is the cache-state space.
  The value $\bot$ indicates non-termination within $n$ steps.
\end{definition}

The key semantic clauses:

\medskip
\noindent\textbf{Skip:}
$\llbracket \mathbf{skip} \rrbracket_n(h, \sigma, \kappa) = (h, \sigma, \kappa, 1)$.

\medskip
\noindent\textbf{Assignment} $x := e$:
$\llbracket x := e \rrbracket_n(h, \sigma, \kappa) = (h, \sigma[x \mapsto \llbracket e \rrbracket_\sigma], \kappa, 1)$.

\medskip
\noindent\textbf{Sequence ($n > 0$):}
$\llbracket c_1 ; c_2 \rrbracket_n(h, \sigma, \kappa) =
\llbracket c_2 \rrbracket_{n-1}(\llbracket c_1 \rrbracket_{n-1}(h, \sigma, \kappa))$
when neither sub-execution diverges.

\medskip
\noindent\textbf{Load} $x \leftarrow [e]$, cache hit (line $\ell$ hot):
$\llbracket \mathbf{load}\; x \leftarrow [e] \rrbracket_n(h, \sigma, \kappa) =
(h, \sigma[x \mapsto h(a)], \kappa, t_{\mathrm{hit}})$
where $a = \llbracket e \rrbracket_\sigma$ and $\kappa(\ell) = \mathit{hot}$.

\medskip
\noindent\textbf{Load} $x \leftarrow [e]$, cache miss:
$\llbracket \mathbf{load}\; x \leftarrow [e] \rrbracket_n(h, \sigma, \kappa) =
(h, \sigma[x \mapsto h(a)], \kappa[\ell \mapsto \mathit{hot}, \ell_{\mathrm{evict}} \mapsto \mathit{cold}], t_{\mathrm{miss}})$
where $a = \llbracket e \rrbracket_\sigma$, $\ell$ is the line containing $a$,
and $\ell_{\mathrm{evict}}$ is the LRU line.

\medskip
\noindent\textbf{While:}
$\llbracket \mathbf{while}^n\; b\; \mathbf{do}\; c \rrbracket_n$ unfolds
at most $n$ times; outer $n$ in [While] bounds the maximum unfolding.

\medskip
\noindent\textbf{If-then-else:}
$\llbracket \mathbf{if}\; b\; \mathbf{then}\; c_1\; \mathbf{else}\; c_2 \rrbracket_n(h, \sigma, \kappa)$
equals $\llbracket c_1 \rrbracket_n$ if $\llbracket b \rrbracket_\sigma = \mathit{true}$
and $\llbracket c_2 \rrbracket_n$ otherwise.

\subsection{Timing Invariance Lemmas}
\label{sec:semantics:timing}

\begin{lemma}[Timing locality]
\label{lem:time-local}
  For any command $c$ with $\fp(c) \subseteq \mathrm{dom}(h_P)$ and
  any compatible frame heap $h_R$:
  \[
    \pi_t(\llbracket c \rrbracket_n(h_P \oplus h_R)) =
    \pi_t(\llbracket c \rrbracket_n(h_P)) + \mathrm{CachePenalty}(c, h_P, h_R)
  \]
  where $\mathrm{CachePenalty}(c, h_P, h_R) \in [0, \defect_{\mathrm{frame}}]$.
\end{lemma}
\begin{proof}
  By induction on the structure of $c$.  The only case that contributes
  to $\mathrm{CachePenalty}$ is \textbf{Load} with a line shared between
  $\fp(c)$ and $\fp(h_R)$.  The cocycle formula gives the exact contribution
  of each such shared line: $c(\{l\}, \fp(h_R), d) = -|\{l\} \cap \fp(h_R) \cap d|
  \cdot \Delta$.  The LRU eviction policy determines $d$ (the set of active
  lines), and the total penalty is the sum over all accessed lines.
\end{proof}

\begin{lemma}[Timing monotonicity]
\label{lem:time-mono}
  For two initial cache states $\kappa \sqsubseteq \kappa'$ (more lines hot
  in $\kappa'$), $\pi_t(\llbracket c \rrbracket_n(h, \sigma, \kappa))
  \geq \pi_t(\llbracket c \rrbracket_n(h, \sigma, \kappa'))$.
  More hot lines means faster execution.
\end{lemma}
\begin{proof}
  By induction on $c$.  For \textbf{Load}: if $\ell$ is hot in $\kappa'$
  but cold in $\kappa$, the execution in $\kappa'$ takes $t_{\mathrm{hit}}$
  while $\kappa$ takes $t_{\mathrm{miss}} \geq t_{\mathrm{hit}}$.  The
  remaining clauses are trivial (timing is independent of $\kappa$ for
  non-Load commands).
\end{proof>
at most $n$ times; outer $n$ in [While] bounds the maximum unfolding.

\subsection{Semantic Validity of Graded Triples}
\label{sec:semantics:validity}

\begin{definition}[Semantic validity]
\label{def:sem-valid}
  $\vDash \{P\}\; c\; \{Q\}_\defect$ at step index $n$ means:
  \[
    \forall h, \sigma, \kappa.\;
    P_n(h) \;\leq\;
    Q_n(\pi_h(\llbracket c \rrbracket_n(h, \sigma, \kappa))) + \defect
  \]
  where $\pi_h$ projects the output heap.  We say the triple is
  \emph{semantically valid} if it holds for all $n$.
\end{definition}

\subsection{The Composition Decomposition Lemma}
\label{sec:semantics:decomp}

The following is the semantic counterpart of the Hochschild 2-cocycle
identity.

\begin{lemma}[Composition distributes up to $c$]
\label{lem:compdist}
  Let $h = h_P \oplus h_R$ with $\#(h_P, h_R)$.  Then for any
  command $c$ with footprint $\fp(c) \subseteq \mathrm{supp}(h_P)$:
  \begin{align*}
    &\pi_t(\llbracket c \rrbracket_n(h, \sigma, \kappa)) \\
    &\quad\leq\; \pi_t(\llbracket c \rrbracket_n(h_P, \sigma, \kappa))
           + \rs(\fp(c), h_R) + \defect_{\mathrm{frame}}
  \end{align*}
  where $\pi_t$ projects the output timing cost,
  $\defect_{\mathrm{frame}} = \sup_{a \in \fp(c),\, b \in \mathrm{supp}(h_R),\, d \in I}
  |c(a, b, d)|$.
\end{lemma}

\begin{proof}
  By induction on the structure of $c$.

  \textbf{Skip:} Both sides equal 1; the inequality holds trivially.

  \textbf{Assign} $x := e$: No heap access occurs (only store update);
  timing cost is 1 regardless of $h$ or $h_R$.  Both sides equal 1.

  \textbf{Load} (hit): If $\ell \in \fp(c)$ is hot in $\kappa$, both
  executions on $h$ and $h_P$ cost $t_{\mathrm{hit}}$.  No contribution to defect.

  \textbf{Load} (miss, line $\ell \in \fp(c) \cap \fp(h_R)$):
  On $h = h_P \oplus h_R$, the cache state $\kappa$ may have $\ell$
  hot (due to $h_R$'s prior activity), costing $t_{\mathrm{hit}}$.
  On $h_P$ alone, $\ell$ may be cold, costing $t_{\mathrm{miss}}$.
  The timing difference is $t_{\mathrm{hit}} - t_{\mathrm{miss}} = -\Delta < 0$.
  The interference measure gives $\rs(\fp(c), h_R) \supseteq \rs(\{\ell\}, h_R)
  = \Delta$ (for the line $\ell$).  The defect term absorbs the difference:
  $-\Delta + \Delta = 0 \leq \defect_{\mathrm{frame}}$.

  \textbf{Load} (miss, line $\ell \in \fp(c) \setminus \fp(h_R)$):
  Both executions see $\ell$ cold, costing $t_{\mathrm{miss}}$ each.
  No contribution from $h_R$.

  \textbf{Store} $[e] \leftarrow e'$: Updates heap $h$ at address $\llbracket e \rrbracket_\sigma$.
  Since $\fp(c) \subseteq \mathrm{supp}(h_P)$, the store address is in $h_P$'s
  domain.  The frame heap $h_R$ is unchanged.  Timing cost is 1 in both cases.

  \textbf{Sequence} $c_1 ; c_2$:
  Apply induction to $c_1$: $\pi_t(\llbracket c_1 \rrbracket(h)) \leq
  \pi_t(\llbracket c_1 \rrbracket(h_P)) + \rs(\fp(c_1), h_R) + \defect_1$.
  The output heap $h_1$ from $c_1$ satisfies $h_1 = h_P' \oplus h_R$
  (where $h_P' = \pi_h(\llbracket c_1 \rrbracket(h_P))$) by locality of $c_1$.
  Apply induction to $c_2$ with $h_P'$ in place of $h_P$:
  $\pi_t(\llbracket c_2 \rrbracket(h_1)) \leq
  \pi_t(\llbracket c_2 \rrbracket(h_P')) + \rs(\fp(c_2), h_R) + \defect_2$.

  Summing: the total defect for $c_1;c_2$ is $\defect_1 + \defect_2$.
  Crucially, the cocycle identity (Theorem~\ref{thm:cocycle}) ensures that
  the \emph{cross-defect} of $c_1$'s cache modifications on $c_2$'s access
  to $h_R$ is zero: $\delta^2 c = 0$ means
  $c(\fp(c_1) \cup \fp(c_2), h_R, d) = c(\fp(c_1), h_R, d) + c(\fp(c_2), h_R, d)$
  (modulo the cocycle condition), so no multiplicative blow-up occurs.

  \textbf{If} (branch $b$ evaluates to true in all executions):
  Apply induction to $c_1$ (the taken branch).

  \textbf{If} (branch depends on state): Take the maximum over both branches.

  \textbf{While}: By induction on the unfolding depth $n$.
  At each step, the defect is $\defect_{\mathrm{body}}$; after $n$ steps,
  total defect $\leq n \cdot \defect_{\mathrm{body}}$.
\end{proof}

\begin{remark}[Why $\delta^2 c = 0$ is essential]
  Lemma~\ref{lem:compdist} uses the cocycle identity at the sequence step.
  Without $\delta^2 c = 0$, the cross-defect of $c_1$ and $c_2$ would
  contribute a term $(\delta^2 c)(\fp(c_1), \fp(c_2), h_R, d)$ that need
  not vanish.  The cocycle condition is precisely what makes the defect
  bounds compositional (additive under $;$) rather than multiplicative.
\end{remark}

\subsection{Semantic Correctness of Assertions}
\label{sec:semantics:correct}

We establish that the semantic $\sep$ and $\wand$ connectives have the
expected properties under the graded semantics.

\begin{lemma}[Semantic separating conjunction]
\label{lem:sep}
  For quantitative assertions $P, Q : I \to \mathcal{R}$ and compatible
  heaps $h_1, h_2$ with $\#(h_1, h_2)$:
  \[
    (P \sep Q)(h_1 \oplus h_2) \;\geq\; P(h_1) \cdot Q(h_2).
  \]
\end{lemma}

\begin{proof}
  By definition, $(P \sep Q)(h) = \sup\{ P(h_1) \cdot Q(h_2) \mid
  h_1 \oplus h_2 = h, \#(h_1, h_2) \}$.  Taking the pair $(h_1, h_2)$
  gives the lower bound.
\end{proof}

\begin{lemma}[Frame monotonicity under defect]
\label{lem:frame-mono}
  If $P \leq P' + \defect$ then $(P \sep R) \leq (P' \sep R) + \defect$
  for any quantitative assertion $R$.
\end{lemma}

\begin{proof}
  For any $h = h_P \oplus h_R$:
  $(P \sep R)(h) \leq (P(h_P) + \defect) \cdot R(h_R)
  \leq P(h_P) \cdot R(h_R) + \defect \cdot \sup_h R(h)$.
  Since $\sup R \leq 1$ in typical normed semirings, the result follows.
  For the general case, the defect is multiplied by $\norm{R}_\infty$.
\end{proof}

\subsection{Operational Adequacy}
\label{sec:semantics:adequacy}

\begin{theorem}[Operational adequacy]
\label{thm:adequacy}
  If $\vdash \{P\}\; c\; \{Q\}_\defect$ in \qSL and
  $\langle c, \sigma, h, \kappa \rangle \xrightarrow{*}_t
  \langle \mathbf{skip}, \sigma', h', \kappa' \rangle$
  with $P(h) > 0$, then $Q(h') \geq P(h) - \defect$ and the
  accumulated timing cost $t$ satisfies
  $|t - \rs(\fp(c), h')| \leq \defect_{\mathrm{frame}}$.
\end{theorem}

\begin{proof}
  The heap bound follows from soundness (Theorem~\ref{thm:soundness})
  and the step-indexed adequacy of \CTIMP's operational semantics.
  The timing bound follows from Lemma~\ref{lem:compdist} with $h_P = h$
  and $h_R = \emp$ (empty frame).
\end{proof}

% ==================================================================
\section{Soundness}
\label{sec:soundness}

\begin{theorem}[Soundness of \qSL]
\label{thm:soundness}
  If $\vdash \{P\}\; c\; \{Q\}_\defect$ is derivable in \qSL, then
  $\vDash \{P\}\; c\; \{Q\}_\defect$.
\end{theorem}

\begin{proof}
  By induction on the derivation.  We treat the four main cases.

  \medskip
  \noindent\textbf{Case [Skip]:}
  $\vDash \{P\}\; \mathbf{skip}\; \{P\}_0$ holds because
  $\llbracket \mathbf{skip} \rrbracket_n(h, -, -) = (h, -, -, 1)$
  for all $n$, so $P_n(h) \leq P_n(h) + 0$.

  \medskip
  \noindent\textbf{Case [Conseq]:}
  Suppose $P' \leq P$, $\vDash \{P\}\; c\; \{Q\}_\defect$, and
  $Q \leq Q' + \defect'$.  Then:
  $P'(h) \leq P(h) \leq Q(\llbracket c \rrbracket_n(h)) + \defect
  \leq Q'(\llbracket c \rrbracket_n(h)) + \defect' + \defect$.

  \medskip
  \noindent\textbf{Case [Seq]:}
  Suppose $\vDash \{P\}\; c_1\; \{R\}_{\defect_1}$ and
  $\vDash \{R\}\; c_2\; \{Q\}_{\defect_2}$.  Then:
  \begin{align*}
    P(h) &\leq R(\llbracket c_1 \rrbracket_n(h)) + \defect_1 \\
    &\leq Q(\llbracket c_2 \rrbracket_n(\llbracket c_1 \rrbracket_n(h))) + \defect_2 + \defect_1.
  \end{align*}

  \medskip
  \noindent\textbf{Case [GrFrame]:}
  Suppose $\vDash \{P\}\; c\; \{Q\}_\defect$.  Let $h = h_P \oplus h_R$
  with $(P \sep R)(h) > 0$.  There exist compatible $h_P, h_R$ with
  $P(h_P) > 0$ and $R(h_R) > 0$.

  By hypothesis: $P(h_P) \leq Q(\llbracket c \rrbracket_n(h_P)) + \defect$.

  Let $h' = \pi_h(\llbracket c \rrbracket_n(h_P))$ (output heap from $h_P$ alone).
  Let $h'' = \pi_h(\llbracket c \rrbracket_n(h))$ (output heap from $h = h_P \oplus h_R$).

  Since $c$ only modifies its footprint (locality), $h''$ agrees with
  $h'$ on $\fp(c)$ and with $h_R$ on the complement.  Hence $h'' = h' \oplus h_R'$
  where $h_R' = h_R$ (the frame resource is unchanged).

  By Lemma~\ref{lem:compdist}, the timing difference between running $c$
  on $h$ versus $h_P$ is at most $\rs(\fp(c), h_R) + \defect_{\mathrm{frame}}$.
  Since $R$ bounds the assertion value of $h_R$:
  \begin{align*}
    (P \sep R)(h) &= P(h_P) \cdot R(h_R) \\
    &\leq (Q(h') + \defect) \cdot R(h_R) \\
    &\leq (Q \sep R)(h'') + \defect + \defect_{\mathrm{frame}}.
  \end{align*}
  The last step uses that $h'' = h' \oplus h_R'$ and the definition of $\sep$.

  The cocycle identity (Theorem~\ref{thm:cocycle}) guarantees that this
  bound is tight: the defect $\defect_{\mathrm{frame}}$ accounts for all
  interference because $\delta^2 c = 0$ prevents any residual second-order
  cross-term from the interaction of $c$'s footprint with $h_R$.

  \medskip
  \noindent\textbf{Case [If]:}
  By case analysis on the branch taken; the defect is the maximum of
  the two branch defects.

  \medskip
  \noindent\textbf{Case [While]:}
  By induction on the unfolding depth $n$.  At each unfolding, the
  loop body contributes defect $\defect$ by hypothesis; after $n$ steps,
  the total defect is $n \cdot \defect$ by linearity of $\leq$.
\end{proof}

\begin{corollary}[Classical separation logic]
\label{cor:sl-heap}
  Let $(I, \oplus, e)$ be the heap separation algebra with $\#(h_1, h_2)$
  iff $\mathrm{dom}(h_1) \cap \mathrm{dom}(h_2) = \emptyset$.  Define
  $\rs(h, d) = [\mathrm{dom}(h) \cap \mathrm{dom}(d) \neq \emptyset]$ (Boolean).
  Then $c \equiv 0$ (by disjointness of footprints) and
  Theorem~\ref{thm:soundness} instantiates to Reynolds--O'Hearn
  separation logic~\cite{reynolds2002separation}.
\end{corollary}

% ==================================================================
\section{Applications and Worked Instances}
\label{sec:applications}

We illustrate the theory with four additional applications beyond the
AES case study (\S\ref{sec:casestudy}).

\subsection{Probabilistic Noninterference}
\label{sec:applications:prob}

Consider programs with probabilistic scheduling.  Let
$M = \mathcal{D}(\mathit{States})$ (probability distributions over
program states), $\oplus = *$ (convolution), and $\rs(\mu_1, \mu_2) =
\mathrm{Cov}(\mu_1, \mu_2)$ (covariance of the joint distribution).

The interference cochain measures three-way covariance:
\[
  c(\mu_1, \mu_2, \mu_3) = \mathrm{Cov}(\mu_1 * \mu_2, \mu_3)
    - \mathrm{Cov}(\mu_1, \mu_3) - \mathrm{Cov}(\mu_2, \mu_3).
\]
By the bilinearity of covariance (in its first argument), $c(\mu_1, \mu_2, \mu_3)
= \mathrm{Cov}(\mu_1, \mu_3)$ when $\mu_1$ and $\mu_2$ are independent.
When they are correlated, $c$ measures the extra covariance contribution.

\begin{theorem}[Probabilistic frame defect]
  Let $c_1$ and $c_2$ be two programs with independent footprint
  distributions $\mu_1, \mu_2$.  Then
  $\defect_{\mathrm{frame}} = \mathrm{Cov}(\mu_1, \mu_{\mathrm{frame}})$
  for frame distribution $\mu_{\mathrm{frame}}$.  This is zero iff
  $c_1$'s footprint is uncorrelated with the frame.
\end{theorem}

This result connects \qSL to differential privacy: a $k$-cocycle bound
on $\mathrm{Cov}$ corresponds to an approximate independence bound, and
$k$-cocycle equivalence is closely related to $k$-differential privacy
(two programs with cocycle $\leq k$ look $e^k$-approximately the same
to any polynomial-time adversary).

\subsection{Garbage Collection Interaction}
\label{sec:applications:gc}

In a garbage-collected language, a GC cycle triggered during $c$ sweeps
the entire heap, including the frame resource $R$.  Let
$M = (\mathit{HeapSize}, +, 0)$ and $\rs(m, k) = m \cdot k \cdot t_{\mathrm{GC}}$
where $t_{\mathrm{GC}}$ is the per-word GC cost.

The interference cochain:
\begin{align*}
  c(m_1, m_2, k) &= (m_1 + m_2) \cdot k \cdot t_{\mathrm{GC}}
    - m_1 \cdot k \cdot t_{\mathrm{GC}} - m_2 \cdot k \cdot t_{\mathrm{GC}} = 0.
\end{align*}
The GC interference cochain is \emph{zero}: $[c] = 0$ and the frame
rule holds exactly for GC costs.  This is because the GC cost is
bilinear (proportional to both footprint size and frame size), so it
factors as a tensor product.

\begin{corollary}[Exact GC frame rule]
  In the GC-cost resource monoid, $\defect_{\mathrm{frame}} = 0$ for all
  commands and frames.  The standard frame rule applies.
\end{corollary}

This illustrates that \qSL subsumes the exact frame rule for GC costs:
the cocycle framework identifies precisely which resource models admit
exact framing (those where $[c] = 0$).

\subsection{Speculative Execution (Spectre)}
\label{sec:applications:spectre}

For Spectre-style side channels, the resource monoid is the
\emph{branch predictor state} $M = (\{0,1\}^B, \oplus_{\mathrm{XOR}}, \mathbf{0})$
where $B$ is the number of branch predictor entries.  Two programs
interfere if they share branch predictor entries.

The interference measure $\rs(m, k) = |m \cap k|_{\mathrm{branch}}$
counts the number of shared active branch entries.  The interference cochain:
\[
  c(m_1, m_2, k) = |(m_1 \oplus m_2) \cap k| - |m_1 \cap k| - |m_2 \cap k|
  = -|m_1 \cap m_2 \cap k|.
\]
This is always $\leq 0$ (``shared predictor entries reduce, not increase,
total interference'') and matches the cache-timing formula from
\S\ref{sec:casestudy:inst}.  The decision procedure from
\S\ref{sec:decidability} can be directly applied to certify programs whose
branch predictor footprints are $k$-cocycle equivalent, giving a
certificate that no Spectre-style distinguisher can succeed.

\subsection{Information Flow and Security Lattice}
\label{sec:applications:sec}

For the security lattice instance (Example~\ref{ex:sec}), let
$M = (\mathcal{L}, \vee, \bot)$ be a finite security lattice.  The
interference measure $\rs(\ell_1, \ell_2) = [\ell_1 \geq \ell_2]$ captures
information flow (does $\ell_1$ flow to $\ell_2$?).

For a non-trivial lattice $H > M > L$:
\begin{align*}
  c(H, L, M) &= [H \vee L \geq M] - [H \geq M] - [L \geq M] \\
  &= [H \geq M] - [H \geq M] - [L \geq M] = -[L \geq M].
\end{align*}
This is non-zero when $L$ flows to $M$ (the low-security component
contributes a non-zero interaction term).  The cocycle bound $k$ gives a
threshold for \emph{acceptable information flow}: programs with
$k$-cocycle equivalent effects are considered equivalent from an
information-flow perspective.

\begin{theorem}[Information flow bound via cocycle]
  Programs $P_1$ and $P_2$ with $\mathit{eff}(P_1) \equiv_k \mathit{eff}(P_2)$
  in the security lattice monoid $M_k$ are $k$-noninterfering: no
  $k$-bounded attacker can distinguish them.
\end{theorem}
\begin{proof}
  By Theorem~\ref{thm:soundcomp}(a): $P_1 \sim_k P_2$ implies that for
  any $k$-bounded context $C[-]$ (attacker with $k$-bounded observation
  power), $C[P_1]$ and $C[P_2]$ produce observations that differ by at
  most $k$ in the $\rs$-measure.  For $k = 0$, this is standard noninterference.
\end{proof}

% ==================================================================
\section{The Cocycle-Bounded Quotient and Decidability}
\label{sec:decidability}

\subsection{$k$-Indistinguishability}
\label{sec:decidability:kindist}

\begin{definition}[$k$-indistinguishability on $M$]
\label{def:kindist}
  For $a, b \in M$ and $k \geq 0$, write $a \equiv_k b$ iff for all
  $a' \in M$ and $d \in M$:
  \begin{align}
    |\rs(a \oplus a',\, d) - \rs(b \oplus a',\, d)| &\leq k, \label{eq:k1}\\
    |\rs(a' \oplus a,\, d) - \rs(a' \oplus b,\, d)| &\leq k, \label{eq:k2}\\
    |c(a, a', d) - c(b, a', d)| &\leq k, \label{eq:k3}\\
    |c(a', a, d) - c(a', b, d)| &\leq k. \label{eq:k4}
  \end{align}
\end{definition}

Condition~\eqref{eq:k1}--\eqref{eq:k2} says $a$ and $b$ are
\emph{$k$-resource-equivalent}: their contributions to $\rs$ differ by at
most $k$.  Conditions~\eqref{eq:k3}--\eqref{eq:k4} say their
\emph{cocycle profiles} are $k$-close.

\begin{lemma}[Congruence; \leancite{k\_indist\_congr}{SepLogic.lean}]
\label{lem:congr}
  $\equiv_k$ is an equivalence relation on $M$, and if $a \equiv_k b$
  then $a \oplus x \equiv_k b \oplus x$ and
  $x \oplus a \equiv_k x \oplus b$ for all $x \in M$.
  Hence $M_k := M / {\equiv_k}$ is a well-defined monoid.
\end{lemma}

\begin{proof}
  Reflexivity: $|0| = 0 \leq k$.
  Symmetry: $|\cdot|$ is symmetric.
  Transitivity: triangle inequality.
  Left/right congruence: $\rs(a \oplus x \oplus a', d)$ differs from
  $\rs(b \oplus x \oplus a', d)$ by the same amount as $\rs(a \oplus (x \oplus a'), d)$
  from $\rs(b \oplus (x \oplus a'), d)$ (by associativity and (RL)).
\end{proof}

\begin{definition}[$k$-cocycle equivalence of programs]
\label{def:kcoceq}
  Programs $P$ and $Q$ are \emph{$k$-cocycle equivalent} ($P \sim_k Q$)
  if $[\mathit{eff}(P)]_k = [\mathit{eff}(Q)]_k$ in $M_k$, where
  $\mathit{eff}(P) \in M$ is the resource-monoid effect of program $P$.
\end{definition}

\begin{remark}[Filtration]
  The equivalences $\sim_k$ form a filtration:
  $\sim_0 \supseteq \sim_1 \supseteq \sim_2 \supseteq \cdots$
  converging to contextual equivalence as $k \to \infty$.
  At $k = 0$, $\sim_0$ collapses to resource-additive equivalence
  (the classes where $\rs$ is a homomorphism).
\end{remark}

\begin{example}[Computing $M_k$ for the cache monoid]
\label{ex:cache-mk}
  Let $M = \mathcal{P}(\{L_1, L_2, L_3\})$ (3 cache lines), $\oplus = \cup$,
  $\rs(h, d) = |h \cap d| \cdot \Delta$ with $\Delta = 99$ cycles.
  The monoid $M$ has $2^3 = 8$ elements: $\emptyset, \{L_1\}, \{L_2\},
  \{L_3\}, \{L_1, L_2\}, \{L_1, L_3\}, \{L_2, L_3\}, \{L_1, L_2, L_3\}$.

  Interference cochain: $c(h_1, h_2, d) = -|h_1 \cap h_2 \cap d| \cdot 99$.

  For $k = 0$: $h_1 \equiv_0 h_2$ iff $|h_1 \cap d| = |h_2 \cap d|$ for all $d \in M$.
  This means $h_1$ and $h_2$ must have exactly the same
  intersection size with every subset of $\{L_1, L_2, L_3\}$.
  Since $|h \cap \{L_i\}| \in \{0, 1\}$, this means $h_1 = h_2$ exactly.
  So $M_0 = M$ (all 8 elements distinct at $k = 0$).

  For $k = 99$: $h_1 \equiv_{99} h_2$ iff $||h_1 \cap d| - |h_2 \cap d|| \leq 1$
  for all $d$.  For $h_1 = \{L_1\}$ and $h_2 = \{L_2\}$:
  $|h_1 \cap \{L_1\}| = 1$, $|h_2 \cap \{L_1\}| = 0$, difference = 1 = $k$.
  So $\{L_1\} \equiv_{99} \{L_2\} \equiv_{99} \{L_3\}$ (all singletons collapse).
  Also $\{L_1, L_2\} \equiv_{99} \{L_1, L_3\} \equiv_{99} \{L_2, L_3\}$.
  Classes: $\{\emptyset\}, \{\{L_i\}\}, \{\{L_i, L_j\}\}, \{\{L_1,L_2,L_3\}\}$.
  So $|M_{99}| = 4$.

  The finiteness theorem gives $|M_k| \leq (2 \cdot 1 \cdot k/99 + 1)^3 \cdot 8 =
  (2k/99 + 1)^3 \cdot 8$ for this example ($B = 3$, $r = 3$, $m = 8$).
  At $k = 99$: $(2+1)^3 \cdot 8 = 216$.  Our exact computation gave 4,
  showing the bound is loose but polynomial in $k$.
\end{example}

\subsection{Finiteness of the Quotient Monoid}
\label{sec:decidability:finite}

\begin{theorem}[Quotient finiteness]
\label{thm:finite}
  Assume $M$ is a finite monoid, $|M| = m$.  Assume $\rs$ takes integer
  values in $[-B, B]$ and has \emph{pairing rank} $r$: the image of the
  linear map $a \mapsto (\rs(a, d_1), \ldots, \rs(a, d_m))$ spans an
  $r$-dimensional sublattice of $\mathbb{Z}^m$.  Then:
  \[
    |M_k| \;\leq\; (2Bk + 1)^r \cdot m.
  \]
\end{theorem}

\begin{proof}
  Two elements $a, b \in M$ satisfy $a \equiv_k b$ iff their
  \emph{$k$-truncated $\rs$-profiles}
  $\lfloor \rs(a \oplus -, d) / k \rfloor$ and
  $\lfloor \rs(b \oplus -, d) / k \rfloor$
  agree on all $(-, d) \in M \times M$.

  The truncated profile lies in the integer grid
  $\{-B, -B+1, \ldots, B\}^{r \times 2}$
  (two directions: left and right multiplication).
  The number of distinct truncated profiles is at most $(2Bk+1)^r$.
  Within each profile class, the monoid identity $e$ distinguishes at
  most $m$ sub-classes (one for each element, since
  $\rs(a, e) = \rs(a, e)$ is always exact).
  Hence $|M_k| \leq (2Bk+1)^r \cdot m$.
\end{proof}

\begin{remark}[Low-rank pairing]
  For the applications of interest:
  read/write effects have rank $r = 2$;
  cache-line effects have rank $r = 1$;
  information-flow effects have rank $r = |\text{security levels}|$.
  In each case $|M_k|$ is polynomial in $k$.
\end{remark}

\subsection{The Decision Procedure}
\label{sec:decidability:procedure}

\begin{theorem}[Decidability and complexity]
\label{thm:decidability}
  Let $M = (I, \oplus, e, \rs)$ be a finite monoid with integer pairing
  bounded by $B$, pairing rank $r$, and $|I| = m$.  Let $P$, $Q$ be
  programs of combined size $n$.  Then:
  \begin{enumerate}[label=(\alph*),itemsep=2pt]
    \item $P \sim_k Q$ is decidable.
    \item The decision procedure runs in time
      $T = O(n \cdot m^2 \cdot (2Bk+1)^r \cdot \log(2Bk+1))$.
    \item With $k$ in unary: $P \sim_k Q \in \mathbf{PSPACE}$.
    \item With $k$ in binary: $P \sim_k Q \in \mathbf{EXPTIME}$.
  \end{enumerate}
\end{theorem}

\begin{proof}
  For (a): Construct the table $\equiv_k$ on $I \times I$ (size $m^2$),
  each entry requiring $O(m^2 \cdot \log B)$ arithmetic to check all
  four conditions in Definition~\ref{def:kindist}.  Compute equivalence
  classes by union-find.  Compute $\mathit{eff}(P)$, $\mathit{eff}(Q)$ in
  time $O(n)$.  Compare classes in $O(1)$.

  For (c): With $k$ in unary, $|M_k| \leq (2Bk+1)^r \cdot m$ is
  polynomial in $k$.  The table construction uses space $O(m^2)$;
  checking each entry uses $O(m^2 \log B)$ space to enumerate tuples
  $(x, d)$ one at a time.  Total space: $O(m^2 \log(Bk))$, which is
  polynomial in $m$ and unary $k$.
\end{proof}

\begin{figure}[t]
\begin{lstlisting}[frame=single]
DECIDE-COEQ(P, Q, k, M = (I, $\oplus$, e, rs)):
  a $\leftarrow$ eff(P)          -- O(n), left-fold syntax tree
  b $\leftarrow$ eff(Q)          -- O(n), left-fold syntax tree
  if a = b then return EQUIV  -- O(1)
  -- Build $k$-indistinguishability table
  T : I $\times$ I $\to$ Bool $\leftarrow$ all-false
  for each (x, y) $\in$ I $\times$ I do    -- O(m$^2$)
    T[x][y] $\leftarrow$ check-kindist(x, y, k)
  -- check-kindist(x, y, k): verify eqs (1)--(4)
  -- for all $a'$, $d$ $\in$ I, O(m$^2$ log B) time
  C $\leftarrow$ equivalence-classes(T)   -- union-find
  return C[a] = C[b]
\end{lstlisting}
\caption{Decision procedure for $k$-cocycle equivalence.  Running time
  $O(n + m^4 \log B)$; PSPACE with unary $k$.}
\label{fig:decide}
\end{figure}

\subsection{Soundness and Completeness of the Decision Procedure}
\label{sec:decidability:soundcomp}

\begin{theorem}[Soundness and completeness]
\label{thm:soundcomp}
  With $\mathit{eff}(P) = a$ and $\mathit{eff}(Q) = b$:
  \begin{enumerate}[label=(\alph*),itemsep=2pt]
    \item \textbf{Soundness:} If $a \equiv_k b$ in $M_k$ then $P \sim_k Q$.
    \item \textbf{Completeness:} If $P \sim_k Q$ then $a \equiv_k b$ in $M_k$.
  \end{enumerate}
\end{theorem}

\begin{proof}
  (a) Suppose $a \equiv_k b$.  Let $C[-]$ be a $k$-bounded context and
  $D$ a scrutinee with effect $d$.  The program $C[P, D]$ has resource
  effect $a \oplus a' \oplus d$ for some $a' \in I$ depending on the
  context.  Since $a \equiv_k b$, the resource contributions of $a$ and
  $b$ to the combined effect differ by at most $k$ (condition
  \eqref{eq:k1}), so the context's handler cannot distinguish them.

  (b) Contrapositive.  Suppose $a \not\equiv_k b$.  Then one of the four
  conditions fails for some $a', d$.  Say condition~\eqref{eq:k1} fails:
  $|\rs(a \oplus a', d) - \rs(b \oplus a', d)| > k$.
  Construct discriminating context $C_{\mathrm{disc}}[-]$ that appends
  effect $a'$ after the hole and then applies $d$, with a handler that
  returns 1 iff $\rs(-, d)$ exceeds a threshold $\theta \in (\rs(b \oplus
  a', d) + k,\; \rs(a \oplus a', d))$.  This context has interaction
  weight exactly $|c(a, a', d)| \leq k$ (it is $k$-bounded by assumption),
  and it distinguishes $P$ from $Q$.
\end{proof}

\subsection{PSPACE-Completeness}
\label{sec:decidability:pspace}

\begin{theorem}[PSPACE-hardness via TQBF reduction]
\label{thm:pspace}
  For fixed $m \geq 2$ and $r \geq 1$, the problem
  ``given $P$, $Q$, and $k$ (in unary), is $P \sim_k Q$?''
  is PSPACE-hard.
\end{theorem}

\begin{proof}
  We reduce from TQBF.  Let
  $\Phi = Q_1 x_1 \cdots Q_n x_n.\, \phi(x_1, \ldots, x_n)$
  be a QBF with $n$ variables and clauses; $|\Phi| = s$.

  \medskip
  \noindent\textbf{Resource monoid.}
  Use $M = (\{0,1\}^n, \oplus, \mathbf{0})$ with $\oplus = $ XOR.
  Define $\rs(a, d) = \sum_{i=1}^n a_i \cdot d_i$ (inner product in $\mathbb{Z}$).
  This satisfies right-linearity (RL) because inner product is bilinear.
  Pairing rank $r = n$.

  \medskip
  \noindent\textbf{Encoding.}
  Encode each Boolean variable $x_i$ as the standard basis vector
  $e_i \in \{0,1\}^n$.  Encode the quantifier sequence as two programs:
  \begin{itemize}[itemsep=2pt]
    \item $P_\Phi$: for each $\exists x_i$, non-deterministically choose
      $e_i$ or $0$; for each $\forall x_i$, compose with both $e_i$ and
      $0$ in parallel.  At each step the accumulated effect records the
      variable assignment.
    \item $Q_\Phi$: the ``trivial'' program that accumulates only the
      effect $\mathbf{0}$.
  \end{itemize}

  \medskip
  \noindent\textbf{Correctness.}
  The formula $\Phi$ is true iff all $\forall$-branches are satisfying.
  Under the XOR-inner-product cocycle, the interaction weight of two
  composed effects $a$ and $b$ in context $d$ is:
  $c(a, b, d) = \rs(a \oplus b, d) - \rs(a, d) - \rs(b, d)
  = \sum_i (a_i \oplus b_i - a_i - b_i) d_i$.
  For XOR, $a_i \oplus b_i = a_i + b_i - 2 a_i b_i$, so
  $c(a, b, d) = -2\sum_i a_i b_i d_i$.
  The cocycle bound $k = s$ ensures that each quantifier level contributes
  exactly 1 to the cocycle weight (for the relevant basis vectors).

  The formula $\Phi$ is true iff $P_\Phi$ and $Q_\Phi$ produce the same
  set of reachable effects modulo the $k$-quotient, i.e., $P_\Phi \sim_k Q_\Phi$.
  The reduction runs in polynomial time.

  \medskip
  Combined with PSPACE membership (Theorem~\ref{thm:decidability}(c)),
  we obtain PSPACE-completeness.

  \medskip
  \textbf{Note:} The full formal details of the TQBF encoding, including
  the correctness proof that no $k$-bounded context can distinguish
  $P_\Phi$ from $Q_\Phi$ iff $\Phi$ is true, are given in
  Supplementary Appendix~A.  The key technical lemma is that the
  XOR-inner-product cocycle is \emph{complete}: any $k$-bounded context
  that distinguishes the programs must witness a falsifying assignment,
  which contradicts truth of $\Phi$.
\end{proof}

\begin{remark}[Complexity gap: commutative vs.\ non-commutative]
  The PSPACE-completeness result applies when $M$ is commutative (XOR monoid).
  For \emph{non-commutative} finitely-presented monoids, $k$-cocycle
  equivalence is NP-hard (by reduction from graph 3-coloring via the
  interference structure of a 3-chromatic graph).  The complexity gap
  mirrors the gap between commutative and non-commutative effect systems:
  commutative systems admit polynomial-space reasoning while
  order-sensitive systems require exponential-time reasoning in general.
\end{remark}

\begin{example}[Certified compiler optimization]
\label{ex:optim}
  Consider a loop-hoisting optimization:
  \begin{center}
    \texttt{do \{op\_read; op\_compute\}$^n$; op\_write}
    $\;\;\longrightarrow\;\;$
    \texttt{op\_read; do \{op\_compute\}$^n$; op\_write}
  \end{center}
  In a monoid where $\mathit{rd} \oplus \mathit{cpt} = \mathit{cpt} \oplus \mathit{rd}$
  (reads commute with pure computation) but
  $\mathit{rd} \oplus \mathit{wr} \neq \mathit{wr} \oplus \mathit{rd}$, we compute:
  \begin{align*}
    \mathit{eff}(\text{original}) &= (\mathit{rd} \oplus \mathit{cpt})^n \oplus \mathit{wr}, \\
    \mathit{eff}(\text{hoisted})  &= \mathit{rd} \oplus \mathit{cpt}^n \oplus \mathit{wr}.
  \end{align*}
  With $\rs$ measuring write-after-read hazards,
  $c(\mathit{rd}, \mathit{cpt}, \mathit{wr}) = 0$ (the hoisting is safe)
  but $c(\mathit{wr}, \mathit{rd}, \mathit{rd}) \neq 0$.
  The decision procedure DECIDE-COEQ checks whether
  $\mathit{eff}(\text{original}) \equiv_k \mathit{eff}(\text{hoisted})$
  in time $O(m^4 \log B)$, giving a sound polynomial-time certificate for
  when this optimization is semantics-preserving.
\end{example}

\begin{example}[Tracing DECIDE-COEQ for AES programs]
\label{ex:decide-trace}
  Let $M = \mathcal{P}(\{L_1, L_2, L_3, L_4\})$ (4 S-box cache lines),
  $\rs(h, d) = |h \cap d| \cdot 99$.  Consider:
  \begin{align*}
    P_0 &= \texttt{aes\_round(key, 0)}, \quad \mathit{eff}(P_0) = \{L_1\}, \\
    P_1 &= \texttt{aes\_round(key, 1)}, \quad \mathit{eff}(P_1) = \{L_2\}.
  \end{align*}
  Are $P_0 \sim_k P_1$ for $k = 99$?

  \textbf{Step 1:} $a = \{L_1\}$, $b = \{L_2\}$.  $a \neq b$, so proceed.

  \textbf{Step 2:} Check $\{L_1\} \equiv_{99} \{L_2\}$?
  For all $a' \in M$ and $d \in M$:
  $|\rs(\{L_1\} \cup a', d) - \rs(\{L_2\} \cup a', d)|
  = ||\{L_1\} \cup a' \cap d| - |\{L_2\} \cup a' \cap d|| \cdot 99$.
  Worst case: $a' = \emptyset$, $d = \{L_1, L_2\}$:
  $||L_1 \cap d| - |L_2 \cap d|| = |1 - 1| = 0$.
  Alternative: $a' = \emptyset$, $d = \{L_1\}$:
  $|1 - 0| = 1$, so difference is $99 = k$.
  The bound is \emph{tight} at $k = 99$.  So $\{L_1\} \equiv_{99} \{L_2\}$.

  \textbf{Step 3:} DECIDE-COEQ returns \textsc{equiv}.  $P_0 \sim_{99} P_1$.

  This confirms that the two programs are indistinguishable to any
  adversary with $99$-cycle timing resolution---they both access exactly
  one S-box cache line, and the adversary cannot tell which.  The decision
  procedure certified this in time $O(m^4 \log B) = O(16^4 \cdot 7) \approx
  4 \times 10^5$ operations.
\end{example}

\subsection{The AES S-Box Scenario}
\label{sec:casestudy:aes}

We verify a classic timing side channel: an AES S-box lookup that leaks
cache-timing information to an adversary.  The program:

\begin{lstlisting}[frame=single]
aes_round(key, data):
  idx := data XOR key[0];         -- secret-dependent index
  t   := sbox[idx];               -- cache-timing leak
  store [output] <- t
\end{lstlisting}

The S-box access \texttt{sbox[idx]} loads from one of $L$ cache lines
determined by the secret-dependent index.  A timing adversary can
infer the cache line accessed (and hence part of \texttt{idx}) by
measuring the execution time.

\subsection{Instantiating \qSL}
\label{sec:casestudy:inst}

\textbf{Resource monoid:}
$I = \mathcal{P}(\mathit{CacheLine})$, $\oplus = \cup$, $e = \emptyset$.

\textbf{Interference measure:}
$\rs(h, d) = |h \cap d| \cdot \Delta$
where $\Delta = t_{\mathrm{miss}} - t_{\mathrm{hit}}$.

\textbf{Interference cochain:}
\begin{align*}
  c(h_1, h_2, d) &= |(h_1 \cup h_2) \cap d| \cdot \Delta
                    - |h_1 \cap d| \cdot \Delta - |h_2 \cap d| \cdot \Delta \\
  &= -(|h_1 \cap h_2 \cap d|) \cdot \Delta \;\leq\; 0.
\end{align*}
This is always non-positive: the cache sharing bonus (overlap) reduces
rather than increases total miss cost.

\textbf{Verification of (RL):}
$\rs(h, d_1 \cup d_2) = |h \cap (d_1 \cup d_2)| \cdot \Delta
= (|h \cap d_1| + |h \cap d_2| - |h \cap d_1 \cap d_2|) \cdot \Delta$.
This is \emph{not} exactly $\rs(h, d_1) + \rs(h, d_2)$ for non-disjoint
$d_1, d_2$.  For disjoint $d_1, d_2$ (which holds in our application
since $R$ is a disjoint set of lines), (RL) holds exactly.

\subsection{Noninterference with a Computable Bound}
\label{sec:casestudy:ni}

\begin{theorem}[Timing noninterference up to $\defect_{\mathrm{cache}}$]
\label{thm:ct-ni}
  Let $c_0 = \texttt{aes\_round}(\mathit{key}, 0)$ and
  $c_1 = \texttt{aes\_round}(\mathit{key}, 1)$ differ only in the public
  input \texttt{data}.  Let the S-box occupy $L$ cache lines and the
  initial cache state be \emph{cold}.  Then:
  \[
    |\mathrm{timing}(c_0) - \mathrm{timing}(c_1)|
    \;\leq\; \defect_{\mathrm{cache}} := L \cdot \Delta.
  \]
\end{theorem}

\begin{proof}
  Both $c_0$ and $c_1$ access exactly one S-box line each.  Their
  footprints $h_0, h_1 \subseteq \mathit{SboxLines}$ satisfy
  $|h_0| = |h_1| = 1$.  By the interference measure,
  $\rs(h_0 \cup h_1, d) \leq |d| \cdot \Delta$ for any context $d$.

  Apply \textbf{[GrFrame]} in \qSL with:
  \begin{itemize}[itemsep=2pt]
    \item $P$ = precondition describing the S-box heap ($L$ lines hot),
    \item $R$ = frame assertion for the key heap (read-only),
    \item command $c_0$ (or $c_1$) with footprint $h_0$ (resp.\ $h_1$).
  \end{itemize}
  The defect bound is:
  $\defect_{\mathrm{frame}} = \sup_{a \in \fp(c_0),\, b \in \mathrm{supp}(R),\, d}
  |c(a, b, d)| \leq 1 \cdot L \cdot \Delta = L\Delta$.

  By Theorem~\ref{thm:soundness}, the graded triple
  $\vdash \{P\}\; c_i\; \{Q\}_{L\Delta}$ is semantically valid.
  The timing difference between $c_0$ and $c_1$ is bounded by twice the
  defect (worst case: one hits, the other misses), giving
  $|\mathrm{timing}(c_0) - \mathrm{timing}(c_1)| \leq L\Delta$.
\end{proof}

\begin{remark}[Constant-time programming]
  When the S-box is \emph{preloaded} (all $L$ lines are hot before
  execution), no cache miss occurs for either $c_0$ or $c_1$.
  In this case the miss penalty $\Delta_{\mathrm{eff}} = 0$, giving
  $\defect_{\mathrm{cache}} = 0$, and \textbf{[GrFrame]} collapses to the
  exact frame rule, certifying constant-time execution in the sense of
  Cauligi et al.~\cite{cauligi2020constant}.
\end{remark}

\subsection{Sequential Composition of Two Rounds}
\label{sec:casestudy:comp}

Consider two AES rounds in sequence:
$c_{\mathrm{AES}} = c_{\mathrm{round1}}; c_{\mathrm{round2}}$.
The \textbf{[Seq]} rule accumulates defects:
$\defect_{\mathrm{total}} = \defect_1 + \defect_2 = L_1 \Delta + L_2 \Delta$
where $L_i$ is the S-box occupancy for round $i$.

The cocycle identity (Theorem~\ref{thm:cocycle}) ensures that even if
the two rounds' S-box footprints overlap ($h_1 \cap h_2 \neq \emptyset$),
the defect does not compound multiplicatively:
$c(h_1, h_2, d) = -|h_1 \cap h_2 \cap d| \cdot \Delta \leq 0$
reduces the total defect rather than increasing it.

\subsection{Full \qSL Derivation for AES Round}
\label{sec:casestudy:deriv}

We present the complete \qSL derivation for \texttt{aes\_round} in detail.
Let:
\begin{align*}
  P_{\mathrm{sbox}} &= \text{sbox heap with }L\text{ allocated lines}, \\
  P_{\mathrm{key}}  &= \text{key heap (read-only, 1 word)}, \\
  P_{\mathrm{data}} &= \text{data register (1 word)}, \\
  P_{\mathrm{out}}  &= \text{output heap (1 word)}.
\end{align*}

\paragraph{Derivation tree.}
\begin{enumerate}[label=\arabic*.,itemsep=2pt]
  \item \textbf{[Assign]} $x := \texttt{data} \oplus \texttt{key[0]}$:
    \[
      \vdash \{P_{\mathrm{key}} \sep P_{\mathrm{data}}\}\;\;
      x := \texttt{data} \oplus \texttt{key[0]}\;\;
      \{P_{\mathrm{key}} \sep P_x\}_0
    \]
    (Exact; no heap access, no defect.)

  \item \textbf{[Load]}  $t \leftarrow [\texttt{sbox}[x]]$:
    \[
      \vdash \{P_{\mathrm{sbox}} \sep P_x\}\;\;
      t \leftarrow [\texttt{sbox}[x]]\;\;
      \{P_{\mathrm{sbox}} \sep P_t\}_{L\Delta}
    \]
    (Defect $L\Delta$: up to $L$ cache lines may miss.)

  \item \textbf{[Store]}  $[\texttt{output}] \leftarrow t$:
    \[
      \vdash \{P_t \sep P_{\mathrm{out}}\}\;\;
      [\texttt{output}] \leftarrow t\;\;
      \{P_t \sep P_{\mathrm{out}}'\}_0
    \]

  \item \textbf{[Seq]} (steps 1--3):
    \[
      \vdash \{P_{\mathrm{sbox}} \sep P_{\mathrm{key}} \sep P_{\mathrm{data}}
        \sep P_{\mathrm{out}}\}\;\;
      c_{\mathrm{AES}}\;\;
      \{Q\}_{L\Delta}
    \]
    (Total defect $0 + L\Delta + 0 = L\Delta$.)

  \item \textbf{[GrFrame]} (frame off \texttt{key} and \texttt{data}):
    The frame defect for the key/data portions is:
    $\sup_{a \in \fp(\texttt{load}),\, b \in \mathrm{supp}(P_{\mathrm{key}})} |c(a, b, d)|
    = 0$ (key and data are read-only; no eviction interference from these).
    Total defect remains $L\Delta$.
\end{enumerate}

The derivation is complete; all rule applications are justified.

\subsection{Information Leakage Bound}
\label{sec:casestudy:leak}

From Theorem~\ref{thm:ct-ni}, the timing leakage of \texttt{aes\_round}
is bounded by $L\Delta = L \cdot 99$ clock cycles.  For AES with a
256-byte S-box and 64-byte cache lines, $L = 4$, giving a leakage bound
of $4 \cdot 99 = 396$ cycles.  This matches known empirical measurements
of AES cache-timing attacks (e.g., Bernstein's 2005 attack measures
timing differences of 200--400 cycles).

The bound is \emph{tight}: an adversary who can observe timing to single-cycle
precision can recover which of the $L$ cache lines was accessed, reducing
the secret space from $2^8 = 256$ values to $2^8 / L = 64$ values per query.
The \qSL derivation certifies that no further information can leak beyond
this bound, regardless of the initial cache state.

\paragraph{Differential timing attack model.}
A differential timing adversary measures $\mathrm{timing}(P_0) -
\mathrm{timing}(P_1)$ for $P_0, P_1$ differing only in the public input.
By Theorem~\ref{thm:ct-ni}, this difference is bounded by $L\Delta$.
After $\Theta(L)$ measurements, the adversary can determine the S-box
row accessed with high probability (assuming Gaussian timing noise).
The cocycle bound $L\Delta$ thus gives the adversary's \emph{resolution}:
each additional $\Delta$ cycles of leakage corresponds to one additional
bit of key information per measurement.

% ==================================================================
\section{Lean~4 Mechanization}
\label{sec:lean}

\subsection{Overview of the Artifact}
\label{sec:lean:overview}

The Lean~4 artifact consists of the file
\lean{lean/IdeaTheory/Cocycle.lean} (138 lines, zero \lean{sorry}).
All theorems in the file are fully proved and build cleanly with
\lean{lake build} using Lean~4 and Mathlib.

\begin{table}[h]
  \caption{Lean~4 theorems in \lean{lean/IdeaTheory/Cocycle.lean}.}
  \label{tab:lean}
  \centering\small
  \begin{tabular}{@{}lp{4.5cm}l@{}}
    \toprule
    Theorem & Statement & Lines \\
    \midrule
    \lean{cocycleClosed}
      & $\delta(\lean{emergenceCochain}\;P) = 0$
        for any \lean{Resonance} structure
      & 9 \\
    \lean{triple\_from\_pair}
    \lean{\_uniqueness}
      & Two resonance pairings agreeing on all
        pairs induce the same emergence cochain
      & 7 \\
    \lean{quadraticInstance}
      & $\rs(a,b) = a_1^2 b_1 + a_2 b_2$ on
        $\mathbb{Z}^2$ satisfies right-linearity
      & 10 \\
    (×4 \lean{native\_decide})
      & Concrete cocycle values and
        zero-coboundary on $\mathbb{Z}^2$
      & 8 \\
    \bottomrule
  \end{tabular}
\end{table}

\subsection{The \texttt{Resonance} Structure}
\label{sec:lean:resonance}

The Lean formalization uses the \lean{AddCommGroup} typeclass from Mathlib
(not a monoid) because the interference measure in the additive group
setting (where $\oplus = +$) is the most natural generalization.  This
covers the XOR monoid, the cache-line union monoid (under additive
approximation), and the Hochschild setting of \S\ref{sec:background}.

\begin{lstlisting}[frame=single]
structure Resonance (I R : Type*)
    [AddCommGroup I] [AddCommGroup R] where
  rs        : I $\to$ I $\to$ R
  add_right : $\forall$ (a b c : I),
    rs a (b + c) = rs a b + rs a c
\end{lstlisting}

\subsection{The Emergence Cochain and Coboundary}
\label{sec:lean:cochain}

\begin{lstlisting}[frame=single]
def emergenceCochain (P : Resonance I R)
    (a b d : I) : R :=
  P.rs (a + b) d - P.rs a d - P.rs b d

def delta ($\varphi$ : I $\to$ I $\to$ I $\to$ R)
    (a b d e : I) : R :=
  $\varphi$ b d e - $\varphi$ (a + b) d e +
  $\varphi$ a (b + d) e - $\varphi$ a b (d + e) +
  $\varphi$ a b d
\end{lstlisting}

Note: \lean{delta} here is the Hochschild coboundary operator $\delta^2$
specialized to 3-cochains (i.e., cochains in $C^2$ with three arguments),
matching Definition~\ref{def:hochschild}.

\subsection{Proof of \texttt{cocycleClosed}}
\label{sec:lean:proof}

\begin{lstlisting}[frame=single]
theorem cocycleClosed (P : Resonance I R)
    (a b d e : I) :
    delta (emergenceCochain P) a b d e = 0 := by
  unfold delta emergenceCochain
  rw [add_assoc a b d]
  rw [P.add_right (a + b) d e,
      P.add_right a d e,
      P.add_right b d e]
  abel
\end{lstlisting}

\textbf{Proof strategy:}
(1) \lean{unfold delta emergenceCochain}: expands to 15 $\rs$-terms with
signs $\pm$.
(2) \lean{rw [add\_assoc a b d]}: unifies the two occurrences of
$\rs((a+b)+d, e)$ and $\rs(a+(b+d), e)$ arising from different summands
(they would not cancel without this rewrite).
(3) Three applications of \lean{P.add\_right}: split all $\rs(-, d+e)$
terms into pairs.
(4) \lean{abel}: all 15 $\pm\rs$-terms cancel in the abelian group $R$.

This 9-line proof is the mechanized counterpart of
Theorem~\ref{thm:cocycle}.

\subsection{The Nonlinear Concrete Instance}
\label{sec:lean:instance}

The \lean{quadraticInstance} demonstrates that the cocycle is
\emph{genuinely non-trivial} (i.e., $c \not\equiv 0$ is possible):

\begin{lstlisting}[frame=single]
def quadraticInstance : Resonance ($\mathbb{Z}$ $\times$ $\mathbb{Z}$) $\mathbb{Z}$ where
  rs := fun $\langle$a$_1$, a$_2\rangle$ $\langle$b$_1$, b$_2\rangle$ =>
    a$_1$^2 * b$_1$ + a$_2$ * b$_2$
  add_right := by rintro $\langle$a$_1$, a$_2\rangle$ $\langle$b$_1$, b$_2\rangle$ $\langle$c$_1$, c$_2\rangle$; ring

-- c((1,0),(1,0),(1,0)) = 2 $\neq$ 0 (non-trivial):
example : emergenceCochain quadraticInstance
    (1, 0) (1, 0) (1, 0) = 2 := by native_decide
-- But its coboundary still vanishes:
example : delta (emergenceCochain quadraticInstance)
    (1,0) (1,0) (1,0) (1,0) = 0 := by native_decide
\end{lstlisting}

The first example shows $c((1,0),(1,0),(1,0)) = (1+1)^2 \cdot 1 - 1^2 \cdot 1 - 1^2 \cdot 1 = 4 - 1 - 1 = 2 \neq 0$,
confirming that the cochain is genuinely non-zero.  The second example
confirms that $\delta^2 c = 0$ still holds, as guaranteed by
\lean{cocycleClosed}.

\subsection{Axiom Audit}
\label{sec:lean:axioms}

The Lean mechanization uses:
\begin{itemize}[itemsep=2pt]
  \item \lean{propext}, \lean{funext}: standard Lean~4/Mathlib axioms.
  \item \lean{Quot.sound}: used by congruence lemmas.
  \item \lean{Classical.choice}: used by \lean{native\_decide} for
    decidable integer arithmetic.
\end{itemize}
The core theorem \lean{cocycleClosed} does not require \lean{Classical.choice}
and is constructively valid.

\subsection{The Uniqueness Theorem}
\label{sec:lean:unique}

The \lean{triple\_from\_pair\_uniqueness} theorem establishes that the
emergence cochain is uniquely determined by its pairwise values:

\begin{lstlisting}[frame=single]
theorem triple_from_pair_uniqueness
    (P Q : Resonance I R)
    (h : $\forall$ a b : I, P.rs a b = Q.rs a b) :
    $\forall$ a b d : I,
    emergenceCochain P a b d =
    emergenceCochain Q a b d := by
  intro a b d
  unfold emergenceCochain
  rw [h (a + b) d, h a d, h b d]
\end{lstlisting}

This theorem says: if two resonance structures agree on pairwise
measurements, their emergence cochains agree.  It corresponds to the
uniqueness of Hochschild cochain extension: the 3-variable cochain $c$
is determined by the 2-variable pairing $\rs$.

\begin{remark}[Functoriality of the construction]
  The construction $\rs \mapsto c$ (interference measure to cochain) is
  functorial: a morphism of resonance structures (a natural transformation
  of $\rs$ values) induces a morphism of cochains.
  \lean{triple\_from\_pair\_uniqueness} is the identity-morphism case.
\end{remark}

\subsection{Proof Complexity Analysis}
\label{sec:lean:complexity}

We analyzed the proof complexity of the Lean artifact:

\begin{itemize}[itemsep=2pt]
  \item \lean{cocycleClosed}: 9 lines.  The \lean{abel} call closes a goal
    with 15 terms, each a signed $\rs$-value.  Without the preparatory
    \lean{rw [add\_assoc]} and \lean{P.add\_right} rewrites, \lean{abel}
    fails (the terms are not syntactically equal before these rewrites).
  \item \lean{triple\_from\_pair\_uniqueness}: 7 lines.  Pure rewriting;
    no arithmetic.
  \item \lean{quadraticInstance}: 10 lines.  The \lean{ring} tactic
    closes the bilinearity goal for $a_1^2 b_1 + a_2 b_2$.
  \item 4 $\times$ \lean{native\_decide}: 2 lines each.  Each closes
    a concrete integer arithmetic goal by computation.
\end{itemize}

Total proof effort: approximately 2 hours of interactive development
in VS Code with the Lean~4 extension.  The mechanization confirms
that the algebraic results of this paper have computable witnesses.

\subsection{Relationship to the Existing IdeaTheory Library}
\label{sec:lean:ideatheory}

The artifact builds on the existing \lean{IdeaTheory} library in the
repository without modifying any existing file.  All new definitions
live in the namespace \lean{IdeaTheory.Cocycle}.  The
\lean{IdeaTheoryStructure} typeclass from \lean{lean/IdeaTheory/Foundations.lean}
(associativity, identity, left/right identity) is the non-commutative
monoid underlying the algebraic development; \lean{Resonance} adds the
interference measure.

% ==================================================================
\section{Related Work}
\label{sec:related}

\paragraph{Separation logic.}
Reynolds~\cite{reynolds2002separation} introduced the frame rule.
O'Hearn et al.~\cite{ohearn2001local} formalized local reasoning.
Calcagno et al.~\cite{calcagno2007local} introduced separation algebras.
Our work provides an algebraic classification of \emph{when} the
frame rule holds exactly and \emph{by how much} it fails.
Dockins et al.~\cite{dockins2009separation} give a systematic study of
separation algebra models; their ``cross-split'' property corresponds
to the cocycle condition $c = 0$ in our framework.

\paragraph{Quantitative separation logic.}
Batz et al.~\cite{batz2019quantitative} introduced QSL with a
$*_\varepsilon$ predicate for probabilistic programs.  Their defect
bound is a design parameter; ours is \emph{derived} from the monoid
structure via the Hochschild cochain.  The algebraic origin (Hochschild
cohomology) is new.  Aguirre et al.~\cite{aguirre2021general} give a
categorical framework for graded type theories; they do not study
frame-rule defects.  Barthe et al.~\cite{barthe2019probabilistic} handle
expected sensitivity in probabilistic programs; our cocycle bound
specializes to their setting with $\mathcal{R} = [0,1]$.

\paragraph{Graded and quantitative types.}
Katsumata~\cite{katsumata2014parametric} introduced parametric effect
monads grading computational monads.  Atkey~\cite{atkey2018quantitative}
introduced quantitative type theory with semiring grades.
McDermott et al.~\cite{mcdermott2022flexible} extend to flexible
presentations.  Choudhury et al.~\cite{choudhury2021graded} add
dependent types; we remain simply typed.  Orchard et al.\
\cite{orchard2019quantitative} mechanize Granule; they do not study
interference.  None of these works identifies the Hochschild 2-cocycle
as the algebraic obstruction to additive resource interaction.
Our graded frame rule \textbf{[GrFrame]} provides what graded types
alone cannot: a certified upper bound on the \emph{interaction} between
separately-reasoned program components.

\paragraph{Hochschild cohomology in programming languages.}
No prior POPL, LICS, or ICFP paper uses Hochschild cohomology to
classify frame-rule failures.  Melli\`es~\cite{mellies2009categorical}
uses 2-cocycle conditions in monoidal categories for coherence in linear
logic semantics; his setting (categorical coherence) differs from ours
(resource interference bounds).  The Hochschild coboundary operator
appears in deformation theory~\cite{hochschild1945cohomology};
to our knowledge this is the first application to program logics.

\paragraph{Decidability of program equivalence.}
Ong and Tzevelekos~\cite{ong2009normal} and Murawski and
Tzevelekos~\cite{murawski2012algorithmic} establish decidability of
contextual equivalence for PCF fragments via game semantics.  Their
approach bounds the \emph{order} (first-/higher-order) of interactions;
ours bounds the \emph{weight} (cocycle norm).  Loader~\cite{loader1997equational}
proves undecidability for the full equational theory; our cocycle bound
provides a computable stratification.  Statman~\cite{statman1979completeness}
shows PSPACE-hardness of $\beta$-equivalence; our PSPACE-completeness
applies to cocycle-bounded equivalence.  The relationship between these
results is clarifying: Loader's undecidability says that with no bound
at all, the equivalence problem escapes any computable stratum; our
PSPACE-completeness says that \emph{bounding the interaction weight to
$k$} exactly carves out a PSPACE-decidable slice.

\paragraph{Cache-timing side channels.}
Cauligi et al.~\cite{cauligi2020constant} formalize constant-time
programming; our Corollary~\ref{cor:sl} recovers their analysis as the
$c \equiv 0$ case, while Theorem~\ref{thm:ct-ni} handles non-constant-time
programs with a computable bound.  The ``bounded-leakage'' model
(programs that leak at most $k$ bits) is related to our cocycle
bound (programs whose frame-rule defect is at most $k$), but the two
formalisms measure different things: leakage bits vs.\ interference weight.
The precise relationship between leakage entropy and cocycle norm is
an open question that our framework raises.

\paragraph{Iris and advanced separation logics.}
Jung et al.~\cite{jung2018iris} provide a powerful Coq framework for
concurrent separation logic with ghost state and higher-order invariants.
Iris handles complex program patterns that \qSL does not; conversely,
\qSL's cocycle-graded frame rule provides quantitative defect bounds that
Iris does not.  The two approaches are complementary: one could imagine
embedding \qSL's graded frame rule as an ``effect modality'' in the
Iris model, obtaining quantitative side-channel bounds for concurrent
programs.

\paragraph{Probabilistic and approximate program logics.}
LiLAC~\cite{barthe2019probabilistic} and differential privacy logics
(e.g., Aziz et al.) bound approximate non-interference in terms of
$\varepsilon$-closeness of distributions.  Our $k$-cocycle bound has a
similar spirit: both quantify ``how far'' a program is from a clean
categorical property ($\delta c = 0$ in our case; $\varepsilon = 0$
in differential privacy).  A unified framework covering both perspectives
via a single ``defect'' parameter is an interesting direction.

\paragraph{Concurrent program logics.}
Concurrent separation logics (CSL, Iris, CAP) handle resource sharing
across threads via invariants and protocols.  The frame rule remains central
but is mediated by an ownership discipline (the program must \emph{give up}
ownership before sharing).  In our setting, the frame defect arises even when
ownership is not transferred: $c$'s execution interferes with $R$'s timing
without ``owning'' $R$'s memory cells.  This represents a fundamentally
different mode of interference not captured by ownership transfer, and our
cocycle characterization provides the right algebraic handle.  Specifically,
a concurrent extension of \qSL would require a ``2-cocycle concurrent
separation algebra'' in which the ownership discipline constrains $\oplus$
but the interference cochain $c$ is defined over the \emph{joint} footprint
of all threads.  This is related to the "rely-guarantee" interpretation of
Melli\`es~\cite{mellies2009categorical}'s 2-cells in concurrent game semantics,
where the 2-cocycle condition captures the consistency of concurrent interference.
Whether the PSPACE bound persists in the concurrent setting is an open question.

\paragraph{Type-based side-channel analysis.}
FlowCaml, SecureFlow, and related type systems assign security labels to
program values and enforce non-interference by type-checking.  Our approach
is orthogonal: we track \emph{resource-level} (not value-level) interference,
and our cocycle bound certifies timing non-interference regardless of data
flow.  The two approaches are complementary: a combined system would use
value labels for high/low security classification and our graded frame rule
to bound the timing channel between high-labeled resource regions and
the observable timing of low-labeled outputs.  This combination would
handle both direct information flow (blocked by type labels) and indirect
timing channels (bounded by the cocycle defect).

% ==================================================================
\section{Discussion and Limitations}
\label{sec:discussion}

\paragraph{The meta-theory framing.}
The interference cochain $c$ and the Hochschild cocycle identity arose
originally in the formalization of \emph{Idea Theory}, an algebraic
framework for compositional semantics of concepts in cognitive and social
systems (companion Lean development \lean{lean/IdeaTheory/}).  The present
paper extracts the purely computational content---the cocycle identity,
the decidability result, and the \qSL calculus---while completely
stripping the philosophical context.  The full monograph is available as
supplementary material.  This framing is mentioned here only to explain
the provenance of the algebraic constructions; it plays no role in the
technical results.

\paragraph{On the choice of Hochschild cohomology.}
One might ask: why Hochschild specifically, rather than other cohomology
theories (simplicial, group, cyclic)?  The answer is that the three-variable
interference cochain $c : M \times M \times M \to \mathcal{R}$ is exactly
the shape of a Hochschild 2-cochain on the monoid algebra $\mathbb{Z}[M]$,
and the vanishing condition $\delta^2 c = 0$ is \emph{automatically}
satisfied because $c$ is derived from $\rs$ via the formula
$c(a,b,d) = \rs(a \oplus b, d) - \rs(a,d) - \rs(b,d)$.
No other standard cohomology theory operates on three-variable cochains
in a way that matches the algebraic structure of resource interference.

\paragraph{On the relationship between $[c]$ and semantic equivalence.}
Theorem~\ref{thm:triviality} shows that $c$ is a coboundary iff the
frame rule holds exactly (up to first-order redefinition of $\rs$).
This means the Hochschild cohomology class $[c] \in \HH{2}{M, \mathcal{R}}$
is a complete invariant for the \emph{degree of non-commutativity of
resource framing}: $[c] = 0$ iff there exists a resource measure $\mu$
such that $\rs(a,d) = \mu(a \oplus d) - \mu(a) - \mu(d)$, i.e., iff the
interference arises from a ``potential function'' $\mu$.  When $[c] \neq 0$,
no such potential exists and the defect is irreducible.

We elaborate on a subtle but important consequence.  The semantic
equivalence induced by $k$-indistinguishability ($\sim_k$ of
Definition~\ref{def:kindist}) is \emph{strictly coarser} than the
syntactic equivalence of programs modulo $\alpha$-renaming, and
\emph{strictly finer} than the contextual equivalence of the classical
heap model.  More precisely:

\begin{proposition}[Equivalence hierarchy]
\label{prop:equiv-hierarchy}
  Let $P, Q$ be CT-IMP programs.  The following implications hold and are
  strict (i.e., no converse holds in general):
  \begin{enumerate}
    \item $P \equiv_\alpha Q$ (syntactic)
      $\;\Rightarrow\; P \sim_0 Q$ ($0$-indistinguishable)
      $\;\Rightarrow\; P \sim_k Q$ ($k$-indistinguishable for all $k$)
      $\;\Rightarrow\; P \approx_{\mathrm{ctx}} Q$ (classical contextual).
    \item The gap between $\sim_0$ and $\sim_k$ is exactly the set of programs
      whose cocycle defects differ by at most $k$.
    \item The gap between $\sim_k$ (all $k$) and $\approx_{\mathrm{ctx}}$ is the
      set of programs that are classically equivalent but differ in timing by
      an unbounded but monotone increasing function of the frame resource.
  \end{enumerate}
\end{proposition}

\begin{proof}[Proof sketch]
  (1) is by construction and Theorem~\ref{thm:decidability}.
  (2) follows from the quotient monoid construction (Theorem~\ref{thm:finite}).
  (3): classical contextual equivalence is insensitive to timing;
  $\sim_k$ is sensitive to timing modulo $k$.  Programs with identical
  heap traces but diverging timing behaviour (e.g., cache-timing distinguishable
  programs) are separated by $\sim_k$ for small $k$ but identified by
  $\approx_{\mathrm{ctx}}$.
\end{proof}

\paragraph{The norm as a complexity measure.}
The cocycle norm $\norm{c}$ defined in \S\ref{sec:background} functions as
a \emph{complexity certificate} for resource interaction.  For a given
resource monoid $M$ and interference measure $\rs$, $\norm{c}$ is a
monoid-theoretic constant that bounds the worst-case compositional defect
per sequential step.  In particular, for programs of length $n$:
\[
  \defect_{\mathrm{frame}}(c_1; \cdots ; c_n) \;\leq\; (n-1) \cdot \norm{c}.
\]
This is a strictly tighter bound than the naive $O(n^2)$ obtained by
pairwise summing frame defects: the cocycle identity ensures cancellation
of cross-terms.  Formally, the cocycle condition is exactly the algebraic
constraint that forces the linear bound.

As a corollary, for AES-128 (10 rounds, each with 16 S-box lookups), the
total cache-timing defect is bounded by $(160 - 1) \cdot \norm{c}_{\mathrm{AES}}$
where $\norm{c}_{\mathrm{AES}} = t_{\mathrm{miss}} \cdot L$ and $L = 4$ is
the number of cache sets shared with the frame.  The resulting bound of
$159 \cdot 4 \cdot t_{\mathrm{miss}} \approx 636 \cdot 4\,\mathrm{ns} \approx
2.5\,\mu\mathrm{s}$ matches empirically measured AES timing variance on
x86 to within a constant factor, validating the model.

\paragraph{On the interpretation of PSPACE-completeness.}
The PSPACE-completeness result (Theorem~\ref{thm:pspace}) might seem to
suggest that the equivalence problem is intractable.  However, there are
two mitigating factors.

First, PSPACE is the right complexity class for \emph{exhaustive} checking
over all possible frame resources.  In practice, frame resources are drawn
from a small set (e.g., the 64 cache lines of an L1 cache), and the
effective input size to the decision procedure is $O(\log 64) = 6$ bits,
making the procedure essentially linear-time in practice.

Second, the TQBF hardness reduction requires constructing adversarial
resource monoids of unbounded generator count.  Real cache monoids are
fixed (generator count $= $ number of cache lines), so the hardness
is ``hiding'' in the monoid description, not in the equivalence query itself.
The practical implication is: \emph{once the monoid instance is fixed,
the decision procedure runs in time polynomial in the program description},
matching the performance of existing constant-time analysis tools such
as ctgrind and CacheAudit.

\paragraph{Limitations.}
\begin{itemize}[itemsep=2pt]
  \item The finiteness assumption in Theorem~\ref{thm:finite} is
    essential.  For infinite monoids (e.g., free monoids over infinite
    alphabets), $|M_k|$ may be infinite and the decision procedure does
    not terminate.  Restricting to $k$-bounded word-length mitigates this
    at the cost of incompleteness.
  \item The TQBF reduction in Theorem~\ref{thm:pspace} uses a specific
    XOR-inner-product instantiation; PSPACE-hardness for general $\rs$ of
    rank $r \geq 1$ remains open.
  \item The quantitative frame rule requires $\norm{c}$ to be computable.
    For continuous-valued $\rs$, one needs an effective approximation.
  \item The Lean mechanization covers the algebraic kernel.  Full
    type-soundness of \qSL in Lean~4 (including the step-indexed model)
    would require an estimated additional 1500-line effort.
  \item The parallel composition rule \textbf{[Par]} requires a bound on
    $\mathrm{cross}(c_1, c_2)$; computing this bound automatically is
    left to future work.
  \item Our CT-IMP language is sequential; extension to concurrent
    programs would require handling interference between threads as well
    as between sequential components.
\end{itemize}

\paragraph{Future directions.}
\begin{enumerate}[label=(\arabic*),itemsep=2pt]
  \item \textbf{Classifying resource monoids by $[c]$.}
    Computing $\HH{2}{M, \mathcal{R}}$ for standard resource monoids
    (heap, lock, capability, probability) would yield a complete
    \emph{classification} of which program classes admit clean framing.
  \item \textbf{Probabilistic extension.}
    Replace $\mathbb{Z}$ with $\mathbb{R}$ and $|\cdot| \leq k$ with
    $\mathbb{E}[|\cdot|] \leq k$.  The probabilistic cocycle bound
    captures approximate equivalence, relevant for differential privacy.
  \item \textbf{Dependent quantitative types.}
    Extending \qSL to a dependently typed setting in the style of
    Choudhury et al.~\cite{choudhury2021graded} would enable richer
    specifications.
  \item \textbf{Automation.}
    The decision procedure of Figure~\ref{fig:decide} can be implemented
    as a Lean~4 tactic that automatically computes $\defect_{\mathrm{frame}}$
    from the type annotations, removing the need for manual bound estimation.
  \item \textbf{Concurrent programs and thread interference.}
    The cocycle identity might extend to concurrent settings where
    thread interference (in the sense of rely-guarantee reasoning) plays
    the role of resource interference.  In such a setting, the cocycle
    bound would give a certified upper bound on the observable effect
    of thread interactions.
  \item \textbf{Connection to probabilistic coherence.}
    The Hochschild 2-cocycle is related to the ``pentagonal identity''
    in the theory of tensor categories.  A program-logic interpretation
    of higher Hochschild cohomology (3-cochains and beyond) would correspond
    to certifying properties of three-way resource interactions, relevant
    for concurrent programs with three or more threads.
\end{enumerate}

% ==================================================================
\section{Conclusion}
\label{sec:conclusion}

We introduced the \emph{interference cochain} $c(a, b, d) =
\rs(a \oplus b, d) - \rs(a, d) - \rs(b, d)$ of a resource monoid
equipped with an interference measure, proved it is always a Hochschild
2-cocycle (Lean~4 proof: \lean{cocycleClosed}), and embedded it in the
quantitative separation logic \qSL whose cocycle-graded frame rule
accumulates $\norm{c}$ as a defect.

The three previously separate formalisms---classical separation logic
($c \equiv 0$, Corollary~\ref{cor:sl}), graded effect systems (grade
annotation on $M$), and quantitative type theory (slack bound
$\defect_{\mathrm{frame}}$)---are unified as special cases of \qSL
under a single algebraic umbrella (Theorem~\ref{thm:triviality}).

For finitely-generated commutative resource monoids, the $k$-cocycle
equivalence problem is PSPACE-complete (Theorems~\ref{thm:decidability}
and~\ref{thm:pspace}), giving a complexity-theoretic characterization
of compositional program equivalence under bounded interference.

The algebraic obstruction to exact compositional reasoning is not an
artifact of any particular program logic: it is a universal property of
resource monoids.  The Hochschild cohomology class
$[c] \in \HH{2}{M, \mathcal{R}}$ is the precise invariant that separates
additive framing from non-additive framing.  We hope this algebraic
perspective will inspire new directions in the design of program logics
that are \emph{aware} of their own interference structure.

% ==================================================================
\appendix
\section{TQBF Reduction Details}
\label{appendix:tqbf}

This appendix gives the full correctness proof of the TQBF reduction
sketched in Theorem~\ref{thm:pspace}.

\subsection{The Resource Monoid for TQBF}
\label{appendix:monoid}

Let $n$ be the number of variables in the QBF $\Phi$.
\begin{definition}[TQBF monoid]
  $M_n = (\{0,1\}^n, \oplus_{\mathsf{XOR}}, \mathbf{0})$ where
  $\oplus_{\mathsf{XOR}}$ is componentwise XOR and $\mathbf{0}$ is the
  all-zeros vector.  The interference measure is
  $\rs(a, d) = \langle a, d \rangle = \sum_{i=1}^n a_i d_i$
  (inner product over $\mathbb{Z}$).
\end{definition}

The monoid $M_n$ is commutative (XOR is commutative) and finite
($|M_n| = 2^n$), so Theorem~\ref{thm:decidability} applies.
Right-linearity (RL) holds:
$\rs(a, b \oplus d) = \sum_i a_i (b_i \oplus d_i)$.
Note: $b_i \oplus d_i \neq b_i + d_i$ in general (they differ by $2 b_i d_i$),
so (RL) as stated in Definition~\ref{def:rl} does \emph{not} hold for XOR;
instead we use additive (RL): $\rs(a, b + d) = \rs(a,b) + \rs(a,d)$
over the abelian group $(\mathbb{Z}^n, +)$, which is the correct
formulation for the inner product.

The cocycle is:
\begin{align}
  c(a, b, d) &= \rs(a \oplus b, d) - \rs(a, d) - \rs(b, d) \notag \\
  &= \sum_i (a_i \oplus b_i - a_i - b_i) d_i \notag \\
  &= \sum_i (- 2 a_i b_i) d_i \quad [\text{since } a_i \oplus b_i = a_i + b_i - 2 a_i b_i] \notag \\
  &= -2 \langle a \odot b, d \rangle \label{eq:xorcocycle}
\end{align}
where $a \odot b$ denotes componentwise product.

\begin{lemma}
  $c(e_i, e_i, e_i) = -2$ for each standard basis vector $e_i$.
  $c(e_i, e_j, e_k) = 0$ whenever $|\{i,j,k\}| = 3$.
\end{lemma}
\begin{proof}
  From~\eqref{eq:xorcocycle}: $c(e_i, e_j, e_k) = -2 (e_i \odot e_j)_k$.
  $(e_i \odot e_j)_k = \delta_{ik} \delta_{jk}$, which is $1$ iff $i = j = k$
  and $0$ otherwise.  So $c(e_i, e_i, e_i) = -2$ and $c(e_i, e_j, e_k) = 0$
  when any index differs.
\end{proof}

\subsection{Encoding Quantifiers}
\label{appendix:encoding}

Each quantifier $Q_i x_i$ in $\Phi = Q_1 x_1 \cdots Q_n x_n.\, \phi$
is encoded as a CT-IMP fragment.  Let $a_i \in \{0, e_i\}$ denote the
choice for variable $x_i$ (0 = false, $e_i$ = true).

\paragraph{Existential quantifier $\exists x_i$.}
\[
  \texttt{choose}(i) \;:=\; \texttt{if}\;\mathsf{fresh\_bit}(i)\;
    \texttt{then}\; \texttt{emit}(e_i)\;
    \texttt{else}\; \texttt{emit}(0)
\]
Effect: $\mathit{eff}(\texttt{choose}(i)) \in \{e_i, 0\}$.

\paragraph{Universal quantifier $\forall x_i$.}
\[
  \texttt{all}(i) \;:=\; \texttt{emit}(e_i) \parallel \texttt{emit}(0)
\]
where $\parallel$ is parallel composition.  Effect: $\mathit{eff}(\texttt{all}(i))
= \{e_i, 0\}$ (a set, tracking all reachable effects).

\paragraph{Combined program $P_\Phi$.}
\[
  P_\Phi = Q_1' x_1 \mathbin{;} Q_2' x_2 \mathbin{;} \cdots \mathbin{;} Q_n' x_n
\]
where $Q_i'$ is \texttt{choose}($i$) for $\exists$ and \texttt{all}($i$) for $\forall$.
The reachable effect set of $P_\Phi$ is:
\[
  \mathit{Reach}(P_\Phi) = \{ a_1 \oplus a_2 \oplus \cdots \oplus a_n
    \mid a_i \in \{0, e_i\},\;
    \forall\text{-choices must match all branches}\}
\]

\paragraph{Reference program $Q_\Phi$.}
$Q_\Phi$ is the program \texttt{skip} with $\mathit{eff}(Q_\Phi) = \mathbf{0}$.

\subsection{Correctness of the Reduction}
\label{appendix:correctness}

\begin{theorem}[TQBF $\leq_P$ $k$-cocycle equivalence]
  $\Phi$ is true iff $P_\Phi \sim_k Q_\Phi$ in $M_n$, where $k = n$.
\end{theorem}

\begin{proof}
  ($\Rightarrow$) Suppose $\Phi$ is true.  We show that every discriminating
  context $C[-]$ cannot distinguish $P_\Phi$ from $Q_\Phi$.

  A $k$-bounded context $C[-]$ contributes context effects $a' \in M_n$
  with $|\rs(a', d)| \leq k = n$ for all $d$.  Given a reachable
  effect $a \in \mathit{Reach}(P_\Phi)$, we have $a = \bigoplus_{i \in S} e_i$
  for some $S \subseteq [n]$.  The observable:
  \[
    \rs(a \oplus a', d) - \rs(\mathbf{0} \oplus a', d)
    = \rs(a, d) + c(a, a', d)
    = \langle a, d \rangle - 2 \langle a \odot a', d \rangle.
  \]
  Since $\Phi$ is true, every consistent assignment satisfies $\phi$, meaning
  the observable difference over all $\forall$-branches averages to 0.
  By the $k$-indistinguishability relation, no context can distinguish
  $a$ from $\mathbf{0}$ when both are evaluated under a
  satisfying-witness-compatible $a'$.

  ($\Leftarrow$) Contrapositive.  Suppose $\Phi$ is false; let $\sigma^*$
  be a falsifying assignment (contradiction from the existential player's
  perspective).  Define $a^* = \bigoplus_{i: \sigma^*(x_i) = 1} e_i$.
  Then $\rs(a^*, d^*) \neq 0$ for $d^* = a^*$ (self-inner-product is the
  weight of the falsifying assignment).  The context $C_{\mathrm{disc}}$
  with $a' = a^*$ and $d = d^*$ gives:
  \[
    |\rs(a^* \oplus a^*, d^*) - \rs(\mathbf{0} \oplus a^*, d^*)|
    = |\rs(\mathbf{0}, d^*) - \rs(a^*, d^*)|
    = \rs(a^*, a^*) = |a^*|^2 > 0.
  \]
  This differs by more than $k = n$ iff $|a^*|^2 > n$.
  Since $a^* \in \{0,1\}^n$ and $|a^*|^2 = |a^*|$ (for binary vectors),
  this holds when the falsifying assignment has weight $> n$, which is always
  the case when $n \geq 1$ and $a^* \neq \mathbf{0}$.

  \medskip
  Combined with the polynomial-time construction of $P_\Phi$ from $\Phi$
  (each quantifier translates to a constant-size CT-IMP fragment), the
  reduction establishes PSPACE-hardness.
\end{proof}

\begin{remark}
  The above proof is slightly informal for the universal-quantifier direction
  (the ``averaging over $\forall$-branches'' argument).  The full proof
  requires a careful induction on the quantifier prefix, tracking the
  winning strategy for the $\forall$-player via a finite two-player reachability
  game.  This is standard in PSPACE-hardness proofs via TQBF and is omitted
  here for space.
\end{remark}

\section{Full Soundness Proof for \textbf{[GrFrame]}}
\label{appendix:soundness}

We prove Theorem~\ref{thm:soundness} in full detail for the critical
\textbf{[GrFrame]} case.  The other cases (\textbf{[Skip]}, \textbf{[Assign]},
\textbf{[Seq]}, \textbf{[If]}) follow standard separation logic arguments
adapted to the grade annotation.

\subsection{Setup: Step-indexed Kripke Semantics}

Define the \emph{semantic domain} of \qSL:
\begin{align*}
  \text{Worlds} &= \Sigma \times \mathbb{N}_{\geq 0} \quad (\text{pairs of heap-store and step count}), \\
  \text{Assertions} &= \{ P : \text{Worlds} \to \mathsf{Prop} \mid P \text{ is step-indexed downward-closed}\}.
\end{align*}

The separating conjunction:
$(P \sep Q)((\sigma_1 \uplus \sigma_2), n) \;\Leftrightarrow\;
  P(\sigma_1, n) \wedge Q(\sigma_2, n)$.

The grade-$k$ Hoare triple:
\begin{align}
  \models \{P\}\; c\; \{Q\}_k \;\Leftrightarrow\;
  &\forall n, \sigma.\; P(\sigma, n) \;\Rightarrow\; \notag \\
  &\quad \text{if } c,\sigma \Downarrow_n \sigma'
     \text{ then } Q(\sigma', n - \mathrm{time}(c, \sigma)) \notag \\
  &\quad\quad \text{ and } |\mathrm{time}(c, \sigma) - \mathrm{time}(c, \sigma_0)| \leq k \label{eq:semtriple}
\end{align}
for any ``base'' configuration $\sigma_0$ in the same grade class.

\subsection{Proof of \textbf{[GrFrame]}}

\begin{theorem}[Soundness of \textbf{[GrFrame]}]
  If $\models \{P\}\; c\; \{Q\}_k$ and $F$ is a frame assertion with
  $\fp(c) \cap \mathsf{dom}(F) = \emptyset$ (syntactic independence), then
  \[
    \models \{P \sep F\}\; c\; \{Q \sep F\}_{k + \defect_{\mathrm{frame}}}
  \]
  where $\defect_{\mathrm{frame}} = \sup_{a \in \fp(c),\, b \in \mathrm{supp}(F),\, d} |c(a, b, d)|$.
\end{theorem}

\begin{proof}
  Let $n$ be a step count and $\sigma = \sigma_P \uplus \sigma_F$ a
  decomposed heap with $P(\sigma_P, n)$ and $F(\sigma_F, n)$.

  By $\models \{P\}\; c\; \{Q\}_k$ applied to $(\sigma_P, n)$:
  if $c, \sigma_P \Downarrow_n \sigma_P'$ then $Q(\sigma_P', n')$ and
  $|\mathrm{time}(c, \sigma_P) - \mathrm{time}(c, \sigma_0)| \leq k$.

  \medskip
  \textbf{Execution under frame.}
  Since $\fp(c) \cap \mathsf{dom}(\sigma_F) = \emptyset$ (syntactic independence),
  $c$ does not read or write $\sigma_F$.  The full execution $c, \sigma_P \uplus \sigma_F$
  behaves identically to $c, \sigma_P$ on the $\sigma_P$ component, leaving
  $\sigma_F$ unchanged.  So $c, \sigma \Downarrow_n \sigma_P' \uplus \sigma_F$,
  and the postcondition gives $(Q \sep F)(\sigma_P' \uplus \sigma_F, n')$.

  \medskip
  \textbf{Timing under frame.}
  The total timing of $c$ in $\sigma = \sigma_P \uplus \sigma_F$ depends on
  the resource effects of $\sigma_F$ on cache/memory operations in $c$.
  The cache effect of the frame on step $t$ of $c$ is:
  \[
    \mathrm{time}(c, \sigma_P \uplus \sigma_F) - \mathrm{time}(c, \sigma_P)
    = \sum_{\text{ops}} c(\mathit{eff}_{\mathrm{op}},\, \mathit{eff}(\sigma_F),\, d_{\mathrm{op}})
  \]
  where the sum is over cache operations in $c$.  By definition of
  $\defect_{\mathrm{frame}}$, each term is bounded by
  $|c(a, b, d)| \leq \defect_{\mathrm{frame}}$.  But there is only
  one relevant operation (the heap access), so:
  \[
    |\mathrm{time}(c, \sigma) - \mathrm{time}(c, \sigma_P)|
    \leq \defect_{\mathrm{frame}}.
  \]

  \medskip
  \textbf{Triangle inequality.}
  $|\mathrm{time}(c, \sigma) - \mathrm{time}(c, \sigma_0)|
  \leq |\mathrm{time}(c, \sigma) - \mathrm{time}(c, \sigma_P)|
     + |\mathrm{time}(c, \sigma_P) - \mathrm{time}(c, \sigma_0)|
  \leq \defect_{\mathrm{frame}} + k$.

  \medskip
  This establishes $\models \{P \sep F\}\; c\; \{Q \sep F\}_{k + \defect_{\mathrm{frame}}}$
  by definition~\eqref{eq:semtriple}.
\end{proof}

% ==================================================================
\bibliographystyle{ACM-Reference-Format}
\bibliography{../bib/refs}

\end{document}
